\chapimg{fig/sh09.jpg}
\IfFileExists{extra/lead09.tex}{\chapterlead{\subhd{After this chapter you can} explain how structure, tillage and mulch change infiltration and evaporation; compute the two-stage evaporation with a cover factor; estimate bulk density and porosity change by tillage; and describe the feedback loops the application closes between soil, crop and atmosphere.}
}{}
\chapter{Soil Structure, Tillage, Mulch and Feedback Loops}\label{chap:09}
\markboth{Soil Structure, Tillage, Mulch and Feedback Loops}{Soil Structure, Tillage, Mulch and Feedback Loops}
\IfFileExists{extra/bg09.tex}{\section{Introduction: the soil surface as a control}
The surface of a soil is where rain is partitioned into runoff and infiltration and where water leaves as evaporation. Both are controlled by structure (aggregates, pores, crust) and by what covers the soil (residue, mulch, canopy). Tillage and mulch are therefore the farmer's two direct tools on the surface balance. This chapter treats them one after the other and then ties them in the feedback loops of the application, where the state of the soil changes the water balance and the water balance changes the soil in return.

\section{Soil structure and porosity}
The total porosity is $\phi=1-\rho_b/\rho_s$ with $\rho_s\approx2.65$ Mg\,m$^{-3}$. Tillage loosens the soil: a topsoil of $\rho_b=1.45$ is brought to 1.20 by ploughing, so $\phi$ rises from 0.453 to 0.547. The extra porosity is macro-porosity: it raises the saturated conductivity and the infiltration, but it also dries faster and settles back under rain and traffic, with a half-life of one season or so.

\begin{equation}\phi=1-\frac{\rho_b}{\rho_s},\qquad \Delta z=z_0\left(\frac{\rho_{b0}}{\rho_{b1}}-1\right).\end{equation}
\par\noindent where $\phi$ is the porosity (m$^3$\,m$^{-3}$), $\rho_b$ the dry bulk density and $\rho_s$ the particle density (Mg\,m$^{-3}$), $\Delta z$ the swelling of the layer (m), $z_0$ its initial thickness and $\rho_{b0},\rho_{b1}$ the bulk densities before and after tillage.



A layer of $z_0=0.25$ m at $\rho_{b0}=1.45$ loosened to $\rho_{b1}=1.20$ swells by $\Delta z=5.2$ cm; the same mass of soil takes more room.

\section{Tillage and the surface balance}
Four effects of tillage matter for the water balance: the roughness stores water in depressions and delays runoff; the loosened layer raises infiltration; the buried residue changes the cover; the broken capillary continuity cuts the supply to the surface and so reduces the evaporation of the second stage. No-till reverses these: the soil is firmer, residue stays on top, infiltration is held by biopores. The application reads the tillage choice as a set of multipliers on $K_s$, on the surface storage and on the evaporation factor.

The depression storage is $S_d\approx0.112\,RR+0.031\,RR^2-0.012\,RR\,s$ (mm), with random roughness $RR$ in cm and slope $s$ in per cent. For a ploughed field with $RR=2.5$ cm and $s=2$: $S_d=0.41$ mm; for a fine seedbed with $RR=0.8$: $S_d=0.09$ mm.

\section{Mulches: types and effects}
Mulches are of three kinds: \emph{organic} (straw, stubble, wood chips), \emph{synthetic} (plastic film, usually black, clear or white) and \emph{mineral} (gravel, sand). Each acts through the same few physical effects: it blocks the vapour path (a resistance in series), reflects or absorbs radiation, and changes the soil temperature. Organic mulches also add organic matter and feed the biology. Plastic films cut the evaporation almost completely but they stop the rain; they must be used with drip irrigation.

In the model, evaporation under a cover is the bare-soil value multiplied by $(1-f_c)$, where $f_c$ is the cover factor of the mulch, and the stage-one limit $U$ is shortened. Chapter~7 gave the two stages; here we apply them.

\section{Feedback loops}
A \emph{feedback loop} is a chain of effects that returns on itself. Three loops are built into the application. (i) \emph{Structure--infiltration}: tillage raises $K_s$, more water infiltrates, the soil is wetter, aggregates slake, the surface seals, $K_s$ falls. (ii) \emph{Mulch--evaporation}: mulch reduces evaporation, the surface stays moist, residue decomposes, the cover thins. (iii) \emph{Water--crop--atmosphere}: stress closes the stomata, transpiration falls, the canopy warms, $ET_0$ demand and the stress rise further. Each loop is a pair of equations coupled by a state variable, and each can be reinforcing (positive) or balancing (negative).

A simple discrete model of the first loop is $K_{n+1}=K_n+\alpha(K_{\max}-K_n)-\beta\,R_n K_n$ with rainfall energy $R_n$. Its fixed point is $K^*=\alpha K_{\max}/(\alpha+\beta R)$. With $\alpha=0.02$, $K_{\max}=40$ mm\,h$^{-1}$, $\beta=0.0005$ and $R=60$: $K^*=16.0$ mm\,h$^{-1}$; with a doubled rainfall energy, $K^*=10.0$.

}{}
\banner{Soil structure and tillage layers}\label{ch:struct}
\section{Pore index and consolidation}
After tillage the relative pore space $p$ (1 is freshly loosened) relaxes to a consolidated value $p_c=0.30$:
\begin{equation}\boxed{p(t)=p_c+(p_0-p_c)e^{-\tau_ct},\qquad p_0=1.00\ (\text{conventional}),\ 0.72\ (\text{reduced}),}\end{equation}
with $t$ in days since tillage; no-till keeps $p=0.42$. The half-life of loose pores is $\ln2/\tau_c$.

\exbox{Consolidation after tillage}{Compute $p(t)$ at 30 and 120 days for conventional and reduced tillage at the four rates $\tau_c=0.001,0.005,0.01,0.02$ d$^{-1}$, with the half-lives.}\sol \givenline{$\tau_c=0.001,0.005,0.01,0.02$ d$^{-1}$}\ans Half-life $=\ln2/\tau_c$: $693$ d at $0.001$, $139$ d at $0.005$, $69$ d at $0.01$ and $35$ d at $0.02$. For conventional tillage at $\tau_c=0.01$ and $t=30$: $p=0.30+0.70e^{-0.3}=0.30+0.70(0.7408)=0.8186$. At $t=120$: $0.30+0.70e^{-1.2}=0.30+0.70(0.3012)=0.511$. Table~\ref{tab:pore} lists all cases.\begin{table}[!htbp]\centering\small\caption{Relative pore space after tillage.}\label{tab:pore}
\begin{tabular}{lrrrrr}\toprule
\hd{$\tau_c$ (d$^{-1}$)} & \hd{Half-life (d)} & \hd{conv. 30 d} & \hd{conv. 120 d} & \hd{reduced 30 d} & \hd{reduced 120 d}\\\midrule
0.001 & 693 & 0.979 & 0.921 & 0.708 & 0.673\\
0.005 & 139 & 0.902 & 0.684 & 0.661 & 0.531\\
0.010 & 69 & 0.819 & 0.511 & 0.611 & 0.427\\
0.020 & 35 & 0.684 & 0.364 & 0.531 & 0.338\\\bottomrule\end{tabular}\end{table}
\textbf{Reading.} With the fastest consolidation (0.02) a conventionally tilled topsoil is back at 0.33 after 120 days, nearly consolidated; at 0.001 it is still 0.89 after the same time. The hydrological effect of tillage in the application's bars is read from the Daisy runs (E5), not from this curve; the curve is a drawing constant. \finans


\section{The retention curve from texture}
The thin section uses $\phi=\mathrm{clip}(0.36+0.0016\,\mathrm{clay}+0.0009\,\mathrm{silt};0.34,0.50)$, $\theta_r=0.01+0.0006\,\mathrm{clay}$ (percentages) and $\theta(h)=\theta_r+(\phi-\theta_r)\Se(h)$.

\exbox{Water contents from texture}{Compute $\phi$, $\theta_r$, the field-capacity content $\theta(100)$, the wilting content $\theta(15849)$ and the plant-available water $\theta(100)-\theta(15849)$ for the ten soils.}\sol  Worked for the loam (S06: 40 per cent sand, 40 silt, 20 clay): $\phi=0.36+0.0016(20)+0.0009(40)=0.36+0.032+0.036=0.428$; $\theta_r=0.01+0.0006(20)=0.022$. With $\alpha=0.036$, $n=1.56$: $\Se(100)=0.5107$$\Se(100)=0.4663$ and $\Se(15849)=0.0286$. So $\theta(100)=0.022+0.406(0.4663)=0.211$ and $\theta(15849)=0.034$; the plant-available water is $0.178$ cm$^3$ cm$^{-3}$, that is 1.78 mm per cm. Table~\ref{tab:thetaten} gives all ten soils.\begin{table}[!htbp]\centering\small\caption{Water contents from texture (percentages of sand, silt, clay).}\label{tab:thetaten}
\begin{tabular}{lrrrrrrrrr}\toprule
\hd{ID} & \hd{Soil} & \hd{Sand} & \hd{Silt} & \hd{Clay} & \hd{$\phi$} & \hd{$\theta_r$} & \hd{$\theta(100)$} & \hd{$\theta(15849)$} & \hd{AWC$_{\rm vG}$}\\\midrule
S01 & Coarse loam & 45 & 35 & 20 & 0.423 & 0.022 & 0.200 & 0.052 & 0.148\\
S02 & Coarse sand & 88 & 8 & 4 & 0.374 & 0.012 & 0.141 & 0.019 & 0.122\\
S03 & Sand & 92 & 5 & 3 & 0.369 & 0.012 & 0.016 & 0.012 & 0.004\\
S04 & Loamy sand & 82 & 12 & 6 & 0.380 & 0.014 & 0.028 & 0.014 & 0.015\\
S05 & Sandy loam & 65 & 25 & 10 & 0.399 & 0.016 & 0.079 & 0.017 & 0.062\\
S06 & Loam & 40 & 40 & 20 & 0.428 & 0.022 & 0.211 & 0.034 & 0.178\\
S07 & Silt loam & 20 & 65 & 15 & 0.443 & 0.019 & 0.309 & 0.059 & 0.251\\
S08 & Sandy clay loam & 55 & 15 & 30 & 0.421 & 0.028 & 0.192 & 0.043 & 0.149\\
S09 & Clay loam & 30 & 35 & 35 & 0.448 & 0.031 & 0.345 & 0.102 & 0.243\\
S10 & Silty clay loam & 10 & 55 & 35 & 0.465 & 0.031 & 0.413 & 0.166 & 0.246\\\bottomrule\end{tabular}\end{table}
\textbf{Reading.} Between 100 and 15,849 cm the retention curve gives $0.004$ cm$^3$ cm$^{-3}$ for the sand (0.04 mm per cm) rising to $0.25$ for the silt loam and silty clay loam (2.5 mm per cm). The engineered per-texture constants of Chapter~\ref{ch:e1} are 0.6 mm per cm for the sand and 1.7 to 2.0 for the fine soils. The two agree in order of magnitude for the loams, but not for the sands: with $n=2.68$ almost all of the sand's water leaves the pores before 100 cm. The application labels the AWC constants as typical texture-class values, not derived from the study soil, and this comparison shows why the label matters. \finans


\exbox{Jurin's law and the pore picture}{A pore of radius $r$ fills with water at suction $h$ if $r\le r_c(h)$, where Jurin's law gives $r_c=\dfrac{2\sigma\cos\gamma}{\rho_wgh}\approx\dfrac{0.149\ \mathrm{cm}^2}{h}$. (a) Verify the constant $0.149$ cm$^2$. (b) Find the largest water-filled pore at 10, 31.6, 100, 300, 1000 and 15,849 cm. (c) At what suction does a pore of radius 25 $\mu$m empty?}\sol  (a) $\sigma=0.0728$ N m$^{-1}$, $\gamma=0$, $\rho_w=998$ kg m$^{-3}$, $g=9.81$ m s$^{-2}$: $2\sigma/(\rho_wg)=0.1456/(9790.4)=1.487\times10^{-5}$ m$^2$ $=0.1487$ cm$^2$ $\approx0.149$ cm$^2$. (b) $r_c=0.149/h$ cm, that is $1490/h$ $\mu$m, so the radii are $149$, $47.2$, $14.9$, $4.97$, $1.49$ and $0.094$ $\mu$m (Table~\ref{tab:jur} gives the same in $\mu$m). (c) $h=1490/25=59.6$ cm. The rule of the application, ``$r$ in micrometres empties at about 1490 divided by $r$ centimetres'', is the same law. The ``wide pores'' of the hidden-question panel, those that empty at less than 31.6 cm, have $r>47$ $\mu$m.\begin{table}[!htbp]\centering\small\caption{Capillary radius.}\label{tab:jur}
\begin{tabular}{lr}\toprule
\hd{$h$ (cm)} & \hd{$r_c$ ($\mu$m)}\\\midrule
10 & 149.00\\
31.6 & 47.15\\
100 & 14.90\\
300 & 4.97\\
1000 & 1.49\\
15849 & 0.09\\\bottomrule\end{tabular}\end{table} \finans


\section{Hydraulic drivers: regression on clay}
For each soil the layer shows $\Delta_{\rm till}=D(\mathrm{conv})-D(\mathrm{notill})$ and $\Delta_{\rm cons}$, and fits ordinary least squares on clay content $c$:
\begin{equation}\boxed{\hat\Delta=a+bc,\quad b=\frac{\sum_s(c_s-\bar c)(\Delta_s-\bar\Delta)}{\sum_s(c_s-\bar c)^2},\quad a=\bar\Delta-b\bar c.}\end{equation}
\par\noindent where $\hat\Delta$ is the predicted hydraulic effect of the soil, $c$ the clay fraction, $a$ the intercept, $b$ the slope, $c_s$ and $\Delta_s$ the clay and the effect of soil $s$, and $\bar c$ and $\bar\Delta$ their means over the soils.


\exbox{Regression of a hydraulic effect on clay}{The drainage differences of run set E5 are not part of the repository tables, so the same least-squares recipe is applied here to the archived \emph{drainage share of the memory fate} $\phi_D$ (Chapter~\ref{ch:fate}), of the ten soils, against clay. Fit $\hat\phi_D=a+bc$.}\sol  Clay percentages: 20, 4, 3, 6, 10, 20, 15, 30, 35, 35; mean $\bar c=17.8$. $\bar\phi_D=0.3354$. $\sum(c-\bar c)^2=1367.6$ and $\sum(c-\bar c)(\phi-\bar\phi)=-24.287$. Slope $b=-0.0178$ per percentage point of clay, intercept $a=0.652$, correlation $r=-0.726$, $R^2=0.527$. Table~\ref{tab:ols} lists the fitted values. Each additional percentage point of clay lowers the drainage share of the memory by 0.018, but clay explains only 53 per cent of the variance: the fit is a trend and not a law (the retention parameter $n$ explains the memory far better, Spearman $\rho=0.952$, Chapter~\ref{ch:mem}).\begin{table}[!htbp]\centering\small\caption{Least-squares fit of the drainage share of memory on clay.}\label{tab:ols}
\begin{tabular}{lrrrr}\toprule
\hd{ID} & \hd{Soil} & \hd{Clay (\%)} & \hd{$\phi_D$ (data)} & \hd{fitted}\\\midrule
S01 & Coarse loam & 20 & 0.029 & 0.296\\
S02 & Coarse sand & 4 & 0.265 & 0.580\\
S03 & Sand & 3 & 0.846 & 0.598\\
S04 & Loamy sand & 6 & 0.741 & 0.545\\
S05 & Sandy loam & 10 & 0.643 & 0.474\\
S06 & Loam & 20 & 0.289 & 0.296\\
S07 & Silt loam & 15 & 0.147 & 0.385\\
S08 & Sandy clay loam & 30 & 0.285 & 0.119\\
S09 & Clay loam & 35 & 0.087 & 0.030\\
S10 & Silty clay loam & 35 & 0.022 & 0.030\\\bottomrule\end{tabular}\end{table} \finans


\exbox{Depth of root uptake}{The depth holding a fraction $p$ of the uptake is $z_p=k_p\Delta z$, $k_p=\min\{k:\sum_{i\le k}U_i/\sum U_i\ge p\}$. For the loam season of Chapter~\ref{ch:e2} the transpiration drawn from the first twenty layers (mm) was 20.5, 18.2, 17.2, 16.2, 15.2, 14.0, 13.1, 12.0, 11.2, 10.5, 9.8, 8.9, 7.9, 6.8, 5.7, 5.0, 4.1, 3.4, 3.1, 2.7. Find $z_{0.5}$ and $z_{0.9}$.}\sol  The total uptake over all layers is $229.6$ mm. The cumulative shares reach 50 per cent at layer 8 (cumulative 0.551; layer 7 has 0.498), so $z_{0.5}=8\times2.5=20.0$ cm. They reach 90 per cent at layer 21 (0.906; one layer less gives 0.895), so $z_{0.9}=52.5$ cm. Half of the water is taken from the upper 20 cm although the roots reach 82 cm: the root-density function decreases exponentially with depth.
\banner{Mulch and evaporation layers}\label{ch:mulch}
\section{Mulch against bare soil}
The layer evaluates Engine 1 at bare soil ($\varphi=1$), at your mulch strength and at $\varphi=0.1$. With $E_\varphi,I_\varphi,D_\varphi$ from Equations 5.4, 5.7 and 5.9:
\begin{equation}\boxed{\Delta E=E_{1.0}-E_\varphi,\quad\Delta I=I_{1.0}-I_\varphi,\quad\Delta D=D_\varphi-D_{1.0}.}\end{equation}
\par\noindent where $\Delta E$, $\Delta I$ and $\Delta D$ are the savings of evaporation, irrigation and the change of drainage caused by a cover of fraction $\varphi$ (mm), and the subscripts 1.0 (bare) and $\varphi$ (covered) mark the two runs. \finans


\exbox{What a mulch saves}{For a soil the engine gives bare soil ($\varphi=1$): $E=190$, $I=180$, $D=40$; and with mulch $\varphi=0.5$: $E=140$, $I=148$, $D=62$ mm yr$^{-1}$ (illustration values). Find the three savings and the verdict; then repeat with a wet year ($\rho=1.3$) where $E$ falls from $175$ to $130$, $I$ from $120$ to $112$, and $D$ rises from $150$ to $188$.}\sol \givenline{$\varphi=1$ ; $E=190$ ; $I=180$ ; $D=40$ ; ; $\varphi=0.5$ ; $E=140$ ; $I=148$ ; $D=62$ mm yr$^{-1}$ ; $\rho=1.3$}\ans Normal year: $\Delta E=190-140=50$ mm, $\Delta I=180-148=32$ mm, $\Delta D=62-40=22$ mm. The share of the saved evaporation that becomes less irrigation is $32/50=0.64$ and the share that becomes extra drainage is $22/50=0.44$ (the two exceed 1 together, so part of the drainage is not funded from the evaporation saving alone; the rest is storage change and transpiration, see Equation~(\ref{eq:fate2})). Since $\Delta I>\Delta D$ the verdict is ``dry to normal: most of the saved water becomes less irrigation''. Wet year: $\Delta E=45$, $\Delta I=8$, $\Delta D=38$. Now $\Delta D>\Delta I$ and the verdict reads ``wet conditions: much of the water the mulch saves now leaves as drainage''. \finans


\section{Fate of the saved water}
At each rainfall scale $\rho$ the layer partitions the evaporation saved,
\begin{equation}\label{eq:fate2}f_I=\frac{\max\{\Delta I,0\}}{\max\{\Delta E,0.1\}},\ f_D=\frac{\max\{\Delta D,0\}}{\max\{\Delta E,0.1\}},\ f_T=\frac{\max\{\Delta T,0\}}{\max\{\Delta E,0.1\}},\ f_{\rm rest}=\max\{0,1-f_I-f_D-f_T\},\end{equation}
\par\noindent where $f_I$, $f_D$, $f_T$ and $f_{\rm rest}$ are the shares (0 to 1) of the evaporation saved that appear as less irrigation, less drainage, more transpiration and the rest; the maximum with 0.1 mm avoids division by a small number.

and reports as crossover the first $\rho$ with $\Delta D\ge\Delta I$ and $\Delta D>1$ mm.

\exbox{Partition of the saved evaporation}{A mulch of strength $\varphi=0.5$ gives, at rainfall scales $\rho=0.5,\dots,1.4$, the changes of Table~\ref{tab:fatepart} (illustration values shaped like the study's trend: the irrigation saving falls with rain while drainage rises). Compute the partition of Equation~(\ref{eq:fate2}) and the crossover.}\sol \givenline{$\varphi=0.5$ gives ; $\rho=0.5,\dots,1.4$}\ans At $\rho=1.0$: $f_I=26/50=0.52$, $f_D=13/50=0.26$, $f_T=4/50=0.08$, $f_{\rm rest}=1-0.86=0.14$. The other rows are in the table.\begin{table}[!htbp]\centering\small\caption{Partition of the saved evaporation.}\label{tab:fatepart}
\begin{tabular}{lrrrrrrrr}\toprule
\hd{$\rho$} & \hd{$\Delta E$} & \hd{$\Delta I$} & \hd{$\Delta D$} & \hd{$\Delta T$} & \hd{$f_I$} & \hd{$f_D$} & \hd{$f_T$} & \hd{$f_{\rm rest}$}\\\midrule
0.5 & 38 & 33 & 0.6 & 1 & 0.87 & 0.02 & 0.03 & 0.09\\
0.7 & 44 & 31 & 3 & 2 & 0.70 & 0.07 & 0.05 & 0.18\\
0.9 & 48 & 29 & 8 & 3 & 0.60 & 0.17 & 0.06 & 0.17\\
1.0 & 50 & 26 & 13 & 4 & 0.52 & 0.26 & 0.08 & 0.14\\
1.1 & 50 & 21 & 17 & 4 & 0.42 & 0.34 & 0.08 & 0.16\\
1.2 & 49 & 15 & 24 & 4 & 0.31 & 0.49 & 0.08 & 0.12\\
1.3 & 48 & 9 & 29 & 5 & 0.19 & 0.60 & 0.10 & 0.10\\
1.4 & 46 & 5 & 33 & 5 & 0.11 & 0.72 & 0.11 & 0.07\\\bottomrule\end{tabular}\end{table} \finans

The first $\rho$ in the table with $\Delta D\ge\Delta I$ and $\Delta D>1$ is $\rho=1.2$. Between the grid points $\Delta D-\Delta I$ is $-4$ at $\rho=1.1$ and $+9$ at $\rho=1.2$, so linear interpolation puts the crossover at $\rho=1.1+0.1\times4/13=1.131$. The study reports about 60 per cent of the saved water going to reduced irrigation in dry years, about 63 per cent to extra drainage in wet years, and a crossover near $\rho\approx1.2$; the illustration table, which is not the study's, gives $f_I=0.87$ at $\rho=0.5$ and $f_D=0.72$ at $\rho=1.4$, which have the same trend.


\section{The regime map}
For each cell of (mulch factor, rainfall scale) the engine forms the regime index $\mathcal I=(E-D)/\max\{1,E+D\}$: evaporation-dominated if $\mathcal I>0.05$, drainage-dominated if $\mathcal I<-0.05$, on the boundary otherwise; $\mathcal I=+1$ is all evaporation and $-1$ all drainage.

\exbox{Regime index}{Classify the seven cells of Table~\ref{tab:regime} (illustration values of annual $E$ and $D$).}\sol  Compute $\mathcal I=(E-D)/\max\{1,E+D\}$. For the first cell: $(190-12)/202=+0.881$, evaporation-dominated. For the cell $(\varphi=0.5,\rho=1.0)$: $(140-62)/202=+0.386$. For $(0.5,1.3)$: $(130-188)/318=-0.182$, drainage-dominated. For $(0.2,1.0)$: $(70-95)/165=-0.152$.\begin{table}[!htbp]\centering\small\caption{Regime index.}\label{tab:regime}
\begin{tabular}{lrrrrr}\toprule
\hd{$\varphi$} & \hd{$\rho$} & \hd{$E$} & \hd{$D$} & \hd{$\mathcal I$} & \hd{Regime}\\\midrule
1.0 & 0.7 & 190 & 12 & +0.881 & evaporation\\
1.0 & 1.0 & 185 & 45 & +0.609 & evaporation\\
1.0 & 1.3 & 175 & 150 & +0.077 & evaporation\\
0.5 & 0.7 & 125 & 15 & +0.786 & evaporation\\
0.5 & 1.0 & 140 & 62 & +0.386 & evaporation\\
0.5 & 1.3 & 130 & 188 & -0.182 & drainage\\
0.2 & 1.0 & 70 & 95 & -0.152 & drainage\\\bottomrule\end{tabular}\end{table}
\textbf{Reading.} Strong mulch in a wet climate pushes the soil into the drainage regime where mulch ``saves less'' water for the crop, because the saved evaporation leaves as drainage. \finans


\section{Two-stage evaporation}
With $e_0$ the atmospheric demand under mulch and canopy, $Q_1$ the stage-1 storage and $\tau$ the stage-2 time scale, $t_1=(Q_1-q_0)/e_0$ and
\begin{equation}Q(t)=\begin{cases}q_0+e_0t,&t\le t_1\\ Q_1+2e_0\tau\bigl(\sqrt{1+(t-t_1)/\tau}-1\bigr),&t>t_1\end{cases}\qquad e(t)=\begin{cases}e_0,&t\le t_1\\ e_0/\sqrt{1+(t-t_1)/\tau},&t>t_1\end{cases}\end{equation}
\par\noindent where $Q(t)$ is the cumulative evaporation (mm) and $e(t)$ the rate (mm\,d$^{-1}$) at time $t$ (d), $q_0$ the initial amount, $e_0$ the stage-one rate, $t_1$ the time at which stage one ends, $Q_1$ the amount evaporated by then and $\tau$ the time scale of stage two (d).

The parameters are
$h_i=\frac{0.25}{\alpha}(1.5-0.5\rho)$, $h_s=4/\alpha$, $Q_1=\max\{2,[\theta(h_i)-\theta(h_s)]L_1(0.6+0.4\rho)\}$ with $L_1=40$ mm, $\tau=0.006+0.03(\mathrm{clay}+\mathrm{silt})/100$ d, $q_0=(1-w_f)0.95Q_1$, $w_f=\mathrm{clip}(1-z_d/50;0.05,1)$, and $e_0=e_{0,b}\varphi(1-0.4c)$.

\exbox{Stage-1 and stage-2 evaporation of the ten soils}{With bare soil ($\varphi=1$, $c=0$, so $e_0=4.2$ mm d$^{-1}$), rainfall scale $\rho=1$ and a surface method ($q_0=0$), compute for the ten soils $h_i$, $h_s$, $Q_1$, the stage-1 duration $t_1$, $\tau$ and the cumulative loss at 10 days.}
\textbf{Solution, loam S06 in full.} $\alpha=0.036$: $h_i=0.25/0.036\times(1.5-0.5)=6.94$ cm; $h_s=4/0.036=111$ cm. The retention uses $\phi=0.428$, $\theta_r=0.022$ (Chapter~\ref{ch:struct}): $\Se(6.94)=0.9959$$\Se(6.94)=0.9617$ and $\Se(111)=0.4427$, so $\theta(h_i)-\theta(h_s)=0.406(0.9617-0.4427)=0.2107$ and $Q_1=8.43$ mm. Stage 1 lasts $t_1=Q_1/e_0=2.01$ d. Then $\tau=0.006+0.03(0.60)=0.024$ d. After 10 days: $t-t_1=7.99$ d, $Q=Q_1+2(4.2)(0.024)(\sqrt{1+333.0}-1)=11.9$ mm. Table~\ref{tab:stage} gives all ten soils.\begin{table}[!htbp]\centering\small\caption{Stage-1 storage, duration and cumulative loss in 10 days (bare, $\rho=1$).}\label{tab:stage}
\begin{tabular}{lrrrrrrr}\toprule
\hd{ID} & \hd{Soil} & \hd{$h_i$} & \hd{$h_s$} & \hd{$Q_1$ (mm)} & \hd{$t_1$ (d)} & \hd{$\tau$ (d)} & \hd{$Q(10)$ (mm)}\\\midrule
S01 & Coarse loam & 2.6 & 42 & 5.99 & 1.43 & 0.0225 & 9.5\\
S02 & Coarse sand & 4.7 & 75 & 7.85 & 1.87 & 0.0096 & 10.1\\
S03 & Sand & 1.7 & 28 & 12.71 & 3.03 & 0.0084 & 14.7\\
S04 & Loamy sand & 2.0 & 32 & 11.90 & 2.83 & 0.0114 & 14.2\\
S05 & Sandy loam & 3.3 & 53 & 10.49 & 2.50 & 0.0165 & 13.3\\
S06 & Loam & 6.9 & 111 & 8.43 & 2.01 & 0.0240 & 11.9\\
S07 & Silt loam & 12.5 & 200 & 7.07 & 1.68 & 0.0300 & 11.0\\
S08 & Sandy clay loam & 4.2 & 68 & 7.35 & 1.75 & 0.0195 & 10.6\\
S09 & Clay loam & 13.2 & 211 & 5.62 & 1.34 & 0.0270 & 9.5\\
S10 & Silty clay loam & 25.0 & 400 & 4.60 & 1.10 & 0.0330 & 8.9\\\bottomrule\end{tabular}\end{table}
\textbf{Check against the application.} It says stage 1 lasts about 3 days for the sand at the season-mean demand and about 1 day for the silty clay loam: the table gives $3.03$ and $1.10$ days. Stage 1 is longest in the sands (large $Q_1$, because the window between $h_i$ and $h_s$ scales with $1/\alpha$) while in stage 2 the rate $e_0/\sqrt{1+(t-t_1)/\tau}$ falls faster for the small-$\tau$ sands than for the fine soils: the sands lose their larger stage-1 reserve first and then dry out quickly, as the application describes.


\exbox{Mulch and canopy in the evaporation model}{For the loam, with the stage constants above, compute $e_0$, $t_1$, and the cumulative loss at 10 and 24 days for mulch strengths $\varphi=1,0.5,0.2$ with canopy cover $c=0$ and $c=0.8$.}\sol \givenline{$\varphi=1,0.5,0.2$ ; $c=0$ ; $c=0.8$}\ans $e_0=4.2\varphi(1-0.4c)$. For $\varphi=0.5$, $c=0.8$: $e_0=4.2\times0.5\times0.68=1.428$ mm d$^{-1}$, $t_1=8.43/1.428=5.90$ d. Table~\ref{tab:mulchev} lists the six cases. Mulch lengthens stage 1 in proportion to $1/\varphi$ (from 2.0 days to 10 days for $\varphi=0.2$), and the cumulative loss falls more than in proportion because stage 2 starts later, which is why the evaporation saving grows convexly with mulch strength.\begin{table}[!htbp]\centering\small\caption{Effect of mulch and canopy on evaporation of the loam.}\label{tab:mulchev}
\begin{tabular}{lrrrrr}\toprule
\hd{$\varphi$} & \hd{$c$} & \hd{$e_0$} & \hd{$t_1$ (d)} & \hd{$Q(10)$} & \hd{$Q(24)$}\\\midrule
1.0 & 0 & 4.200 & 2.01 & 11.9 & 14.3\\
1.0 & 0.8 & 2.856 & 2.95 & 10.6 & 12.4\\
0.5 & 0 & 2.100 & 4.01 & 9.9 & 11.2\\
0.5 & 0.8 & 1.428 & 5.90 & 9.3 & 10.2\\
0.2 & 0 & 0.840 & 10.03 & 8.4 & 9.4\\
0.2 & 0.8 & 0.571 & 14.76 & 5.7 & 8.9\\\bottomrule\end{tabular}\end{table} \finans


\exbox{The drying front}{After stage 1, the depth of the dry layer is $\ell(t)=\min\{60,(Q(t)-Q_1)/\Delta\theta\}$ mm with $\Delta\theta=\mathrm{clip}(\theta(h_i)-\theta(10^4);0.05,0.4)$. For the loam bare soil at 10 days, find $\ell$ and the suction at depth $z=40$ mm from $\log_{10}h(z)=\log_{10}h_s+(\log_{10}h_i-\log_{10}h_s)(1-e^{-(z-\ell)/\lambda})$, $\lambda=L_1/3$.}\sol \givenline{$\ell(t)=\min\{60,(Q(t)-Q_1)/\Delta\theta\}$ mm ; $z=40$ mm ; $\lambda=L_1/3$}\ans $\Delta\theta=\theta(6.94)-\theta(10^4)=0.375$. At 10 days $Q=11.9$ mm, so $\ell=(11.9-8.43)/0.375=9.3$ mm. With $\lambda=13.3$ mm, at $z=40$ mm the exponent is $-(40-9.3)/13.3=-2.30$ and $\log_{10}h=0.962$, that is $h=9$ cm.
\banner{Feedback loops and the hidden questions of the vadose zone}\label{ch:fb}
\section{The seven effects}
For each control the engine moves one step on its list and re-evaluates Engine 1; the change is printed and the arrow width is $w_k=5+\mathrm{clip}(\mu_k;0,4)$ pixels, with $\mu_k$ from the table: temperature $\mu_1=8|\Delta Y|$; rain $\mu_2=|\Delta D|/12$; mulch $\mu_3=|\Delta E|/14$; trigger $\mu_4=|\Delta I|/8$; nitrate cost of drainage $\mu_5=|\Delta D|/10$; nitrogen $\mu_6=|\Delta N_L|/8$; consolidation $\mu_7=|\Delta D|/2$. \finans

\exbox{Arrow widths}{One-step differences for a soil (illustration values: $\Delta Y=-0.11$ Mg ha$^{-1}$, $\Delta D=+9.0$ mm for the rain step, $\Delta E=-14$ mm, $\Delta I=-22$ mm, $\Delta N_L=+5.2$ kg ha$^{-1}$, $\Delta D=-0.6$ mm for consolidation) give the arrow widths of Table~\ref{tab:arrow}.}\sol \givenline{$\Delta Y=-0.11$ Mg ha$^{-1}$ ; $\Delta D=+9.0$ mm ; $\Delta E=-14$ mm ; $\Delta I=-22$ mm ; $\Delta N_L=+5.2$ kg ha$^{-1}$ ; $\Delta D=-0.6$ mm}\ans Example, row 4: $\mu_4=22/8=2.75$, so $w_4=5+2.75=7.75$ pixels. Row 3: $\mu_3=14/14=1.0$, $w_3=6.0$. Row 1: $\mu_1=8\times0.11=0.88$, $w_1=5.88$. A step that would give $\mu>4$ is clipped at width 9 pixels.\begin{table}[!htbp]\centering\small\caption{Arrow widths of the seven feedbacks.}\label{tab:arrow}
\begin{tabular}{lrrrrr}\toprule
\hd{\#} & \hd{Control step} & \hd{Effect} & \hd{Value} & \hd{$\mu$} & \hd{Width (px)}\\\midrule
1 & Temperature $+0.5$ K & $\Delta Y$ & -0.11 & 0.880 & 5.88\\
2 & Rain $+0.1$ & $\Delta D$ & +9.00 & 0.750 & 5.75\\
3 & Stronger mulch & $\Delta E$ & -14.00 & 1.000 & 6.00\\
4 & Trigger one step drier & $\Delta I$ & -22.00 & 2.750 & 7.75\\
5 & Drainage (rain step) & $\Delta D$ & +9.00 & 0.900 & 5.90\\
6 & Nitrogen $+15$ per cent & $\Delta N_L$ & +5.20 & 0.650 & 5.65\\
7 & Consolidation one step & $\Delta D$ & -0.60 & 0.300 & 5.30\\\bottomrule\end{tabular}\end{table} \finans


\exbox{The nitrate cost of drainage}{The study reports about $1.06$ kg N ha$^{-1}$ lost per mm of drainage: $\kappa_N=\dfrac{N_L^{(\text{rain step})}-N_L}{D^{(\text{rain step})}-D}$. With the rain step above ($\Delta D=+9.0$ mm, $\Delta N_L=+9.5$ kg N ha$^{-1}$), find $\kappa_N$ and the nitrate lost if the drainage is cut by 40 mm.}\sol \givenline{$\Delta D=+9.0$ mm ; $\Delta N_L=+9.5$ kg N ha$^{-1}$}\ans $\kappa_N=9.5/9.0=1.056$ kg N ha$^{-1}$ mm$^{-1}$, consistent with 1.06. Cutting drainage by 40 mm saves $40\times1.06=42.4$ kg N ha$^{-1}$, equal to a dressing of 42 kg of nitrogen that would otherwise be leached. The row is printed only if $|\Delta D|>0.5$ mm (otherwise ``no change''), because the ratio of two tiny differences is unstable. \finans


\section{Hidden question 1: wetting and drying}
The wetting curve has twice the $\alpha$ of the drying curve. A sensor reads a suction from the water content; the same water content therefore corresponds to different suctions on the two branches.

\exbox{The same water content on two branches}{The loam ($\alpha_d=0.036$, $n=1.56$) is at the trigger on the drying curve, $h_d=300$ cm. Find the suction that the same water content has on the wetting curve ($\alpha_w=2\alpha_d$).}\sol \givenline{$\alpha_d=0.036$ ; $n=1.56$ ; $h_d=300$ cm ; $\alpha_w=2\alpha_d$}\ans $\Se=\Se(300)=0.2615$ on the drying curve. On the wetting branch, $h_w=\alpha_w^{-1}(\Se^{-1/m}-1)^{1/n}$; since $\alpha_w=2\alpha_d$, $h_w=h_d/2=150$ cm exactly (check: $150.0$ cm). A sensor that has just been wetted thus reads 150 cm where a drying soil at the same water content reads 300 cm: the trigger set in suction fires at twice the water deficit after wetting as after drying. Hysteresis makes the same reading mean different amounts of water. \finans


\section{Hidden question 2: where does water still move}
The capillary conductivity is $K_s K_r(h)$, and a film term proportional to $h^{-1.5}$ and anchored at $10^{-7}$ cm d$^{-1}$ at pF 4.2 ($h=15\,849$ cm) is added.

\exbox{Capillary and film conductivity}{For the loam with $K_s=24.96$ cm d$^{-1}$ compute the capillary conductivity, the film term and their sum at suctions 300, 3000, 15,849, $10^5$ and 316,228 cm (pF 5.5), and find the suction at which the film term overtakes the capillary term.}\sol \givenline{$K_s=24.96$ cm d$^{-1}$ compute}\ans At pF 4.2 the film term is $10^{-7}$ cm d$^{-1}$ by construction. At pF 5.5: $10^{-7}(316228/15849)^{-1.5}=10^{-7}\times0.01122=1.12\times10^{-9}$ cm d$^{-1}$. The capillary term is $K_sK_r(h)$ with $K_r$ from the Mualem equation.\begin{table}[!htbp]\centering\small\caption{Conductivity (cm d$^{-1}$) of the loam, capillary and film terms.}\label{tab:film}
\begin{tabular}{lrrr}\toprule
\hd{$h$ (cm)} & \hd{Capillary} & \hd{Film} & \hd{Sum}\\\midrule
300 & 9.50e-04 & 3.84e-05 & 9.88e-04\\
3{,}000 & 3.92e-07 & 1.21e-06 & 1.61e-06\\
15{,}849 & 1.37e-09 & 1.00e-07 & 1.01e-07\\
100{,}000 & 2.61e-12 & 6.31e-09 & 6.31e-09\\
316{,}228 & 5.20e-14 & 1.12e-09 & 1.12e-09\\\bottomrule\end{tabular}\end{table}
The two terms are equal at $h\approx1{,}653$ cm; beyond it, the film dominates. Where the soil is drier than that the conductivity falls only as $h^{-1.5}$ (a factor of 31.6 per decade of suction) instead of by orders of magnitude, so some water still moves in air-dry soil. The film term is an \emph{illustrative} form, not calibrated. \finans


\section{Hidden question 5: where does the nitrate travel}
The panel shows the pores that empty at less than 31.6 cm (the \emph{wide} pores, radius above $47\ \mu$m) and the share of pore volume and of conductivity that they carry: share of volume $1-\Se(31.6)$; share of conductivity $1-K_r(31.6)$.

\exbox{Wide pores and conductivity}{Compute these two shares for the ten soils.}\sol  For the loam: $\Se(31.6)=0.751$ and $K_r(31.6)=0.032$; the wide pores hold 25 per cent of the pore volume and carry 96.8 per cent of the conductivity.\begin{table}[!htbp]\centering\small\caption{Share of pore volume and of conductivity in pores emptied below 31.6 cm.}\label{tab:wide}
\begin{tabular}{lrrr}\toprule
\hd{ID} & \hd{Soil} & \hd{Volume share} & \hd{Conductivity share}\\\midrule
S01 & Coarse loam & 0.359 & 0.9979\\
S02 & Coarse sand & 0.363 & 0.9874\\
S03 & Sand & 0.923 & 1.0000\\
S04 & Loamy sand & 0.830 & 0.9998\\
S05 & Sandy loam & 0.574 & 0.9957\\
S06 & Loam & 0.249 & 0.9677\\
S07 & Silt loam & 0.115 & 0.9330\\
S08 & Sandy clay loam & 0.335 & 0.9914\\
S09 & Clay loam & 0.093 & 0.9514\\
S10 & Silty clay loam & 0.040 & 0.9321\\\bottomrule\end{tabular}\end{table}
\textbf{Reading.} Where the wide pores carry nearly all of the conductivity, solute moves fast through a small part of the pore volume, so a nitrate pulse reaches the drain earlier than the mean water content suggests. \finans


\section{Hidden question 4: the water table}
The water content above a water table at height $z$ is $\Se=[1+(\alpha z)^n]^{-m}$ (equilibrium, suction $h=z$); the quasi-saturated height is where $\Se<0.9$, that is $z_{90}=\alpha^{-1}(0.9^{-1/m}-1)^{1/n}$.

\exbox{How far the water table reaches into the root zone}{Find $\Se$ at 5, 25, 60 and 150 cm above the water table and the height $z_{90}$ for the ten soils.}\sol  For the loam ($\alpha=0.036$, $n=1.56$, $m=0.359$): $0.9^{-1/m}=0.9^{-2.786}=1.341$, so $z_{90}=(1.341-1)^{0.641}/0.036=(0.341)^{0.641}/0.036=0.5013/0.036=13.9$ cm. Table~\ref{tab:wt} lists the ten soils.\begin{table}[!htbp]\centering\small\caption{Water content above a water table and quasi-saturated height.}\label{tab:wt}
\begin{tabular}{lrrrrrr}\toprule
\hd{ID} & \hd{Soil} & \hd{$\Se$(5)} & \hd{$\Se$(25)} & \hd{$\Se$(60)} & \hd{$\Se$(150)} & \hd{$z_{90}$ (cm)}\\\midrule
S01 & Coarse loam & 0.921 & 0.685 & 0.526 & 0.387 & 6.2\\
S02 & Coarse sand & 0.958 & 0.699 & 0.470 & 0.283 & 9.2\\
S03 & Sand & 0.802 & 0.113 & 0.026 & 0.006 & 3.7\\
S04 & Loamy sand & 0.850 & 0.226 & 0.076 & 0.024 & 4.0\\
S05 & Sandy loam & 0.934 & 0.504 & 0.255 & 0.115 & 6.4\\
S06 & Loam & 0.976 & 0.802 & 0.591 & 0.379 & 13.9\\
S07 & Silt loam & 0.989 & 0.911 & 0.786 & 0.603 & 27.8\\
S08 & Sandy clay loam & 0.952 & 0.718 & 0.520 & 0.347 & 8.9\\
S09 & Clay loam & 0.989 & 0.927 & 0.831 & 0.685 & 33.8\\
S10 & Silty clay loam & 0.995 & 0.969 & 0.923 & 0.834 & 79.7\\\bottomrule\end{tabular}\end{table} \finans


\section{Hidden questions 6 and 7: salinity and the uptake curve}
At steady state the salt entering equals the salt leaving: $I\,EC_{\rm irr}=D\,EC_{\rm dr}$. The root-zone salinity is taken as a quarter of the sum, $EC_{\rm rz}=(EC_{\rm irr}+EC_{\rm dr})/4$, and wheat loses $7.1$ per cent of yield per dS m$^{-1}$ above a threshold of 6 dS m$^{-1}$.

\exbox{How salty can the irrigation water be}{For an irrigation of $I=150$ mm and drainage of $60$, $30$ and $15$ mm, find $EC_{\rm dr}$, $EC_{\rm rz}$ and the relative yield for irrigation water of 2, 4 and 6 dS m$^{-1}$. Then find the largest irrigation-water salinity that keeps the yield at 100 per cent.}\sol \givenline{$I=150$ mm}\ans $EC_{\rm dr}=EC_{\rm irr}I/D$. For $I=150$, $D=60$, $EC_{\rm irr}=2$: $EC_{\rm dr}=5$, $EC_{\rm rz}=(2+5)/4=1.75$ dS m$^{-1}$, below 6, so yield 100 per cent. For $D=15$ and $EC_{\rm irr}=6$: $EC_{\rm dr}=60$, $EC_{\rm rz}=16.5$, loss $0.071(16.5-6)=0.745$, yield 25.5 per cent. See Table~\ref{tab:salt}.\begin{table}[!htbp]\centering\small\caption{Salt balance and wheat yield.}\label{tab:salt}
\begin{tabular}{lrrrrr}\toprule
\hd{$I$} & \hd{$D$} & \hd{$EC_{\rm irr}$} & \hd{$EC_{\rm dr}$} & \hd{$EC_{\rm rz}$} & \hd{Yield (\%)}\\\midrule
150 & 60 & 2 & 5.0 & 1.75 & 100.0\\
150 & 60 & 4 & 10.0 & 3.50 & 100.0\\
150 & 60 & 6 & 15.0 & 5.25 & 100.0\\
150 & 30 & 2 & 10.0 & 3.00 & 100.0\\
150 & 30 & 4 & 20.0 & 6.00 & 100.0\\
150 & 30 & 6 & 30.0 & 9.00 & 78.7\\
150 & 15 & 2 & 20.0 & 5.50 & 100.0\\
150 & 15 & 4 & 40.0 & 11.00 & 64.5\\
150 & 15 & 6 & 60.0 & 16.50 & 25.5\\\bottomrule\end{tabular}\end{table}
\textbf{Threshold.} $EC_{\rm rz}=EC_{\rm irr}(1+I/D)/4\le6$ gives $EC_{\rm irr}\le24/(1+I/D)$: $6.86$ dS m$^{-1}$ for $D=60$ ($I/D=2.5$), $4.0$ for $D=30$ and $2.2$ for $D=15$. Less drainage means less leaching and a lower tolerable salinity: the same irrigation water that is safe in a free-draining soil damages wheat in a tight soil. The steady state, the one-quarter rule and the threshold are stated simple forms ``for teaching and comparison'', not calibrated to a field. \finans


\exbox{Reduction of root uptake with suction}{The uptake curve has breakpoints at 10, 25, 500 and 8,000 cm (a Feddes-type reduction function; here assumed linear in $h$ between breakpoints): uptake is zero at $h\le10$ cm (anaerobic), rises linearly to 1 at 25 cm, stays at 1 until 500 cm and falls linearly to 0 at 8,000 cm. Find the relative uptake at $h=15,\ 300,\ 1000,\ 3000$ and 6000 cm.}\sol \givenline{$h=15,\ 300,\ 1000,\ 3000$}\ans At 15 cm: $(15-10)/(25-10)=0.333$. At 300 cm: 1. At 1000 cm: $1-(1000-500)/7500=0.933$. At 3000: $1-2500/7500=0.667$. At 6000: $1-5500/7500=0.267$. A trigger of $-1000$ cm therefore costs a crop only about 7 per cent of the potential uptake at the moment of the trigger, but the average over the dry-down is lower. In the loam $-1000$ cm is a depletion of 0.76, beyond the $f^*=0.45$. \finans


\section{Hidden question 8: keeping up with the sun}
The topsoil delivers water against the demand until stage 2 begins; mulch reduces the demand and lengthens stage 1 (Chapter~\ref{ch:mulch}). The panel's answer is the stage-1 duration $t_1=Q_1/e_0$.
\exbox{Stage 1 with and without mulch}{For the sand (S03, $Q_1=12.7$ mm) and the silty clay loam (S10, $Q_1=4.6$ mm) find how long the topsoil keeps up with a demand of 4.2 mm d$^{-1}$ bare, and with a mulch $\varphi=0.3$.}\sol \givenline{$Q_1=12.7$ mm ; $Q_1=4.6$ mm ; $\varphi=0.3$}\ans Bare: sand $12.7/4.2=3.0$ d; silty clay loam $4.6/4.2=1.1$ d. With mulch $e_0=4.2\times0.3=1.26$: sand $12.7/1.26=10.1$ d; silty clay loam $4.6/1.26=3.7$ d. Mulch triples the time that the surface keeps pace in both soils; the sand has the longer reserve.
\IfFileExists{extra/app09.tex}{\subsection{The output of the application, panel by panel}
Four layers show the effect of structure, tillage and mulch: the microscope view of the topsoil, the comparison of a mulched and a bare soil, the feedback loops of the whole field and the fate of the water saved by the mulch.

The sections are drawn from the pore-size distribution of the model; the white grains are sand, the blue spaces are the pores holding water and the dark spaces are the air-filled pores.

\appfig{v3_0}{0.96}{Tillage and structure of the topsoil, under the microscope: modelled thin sections of the coarse loam for several tillage systems.}

The bars compare the pore space of the three tillage systems and the change of drainage; consolidation closes the large pores and lowers the drainage.

\appfig{v3_1}{0.96}{Relative pore space after tillage, and the change in drainage at the current settings, coarse loam.}

The two columns show the soil with and without mulch and the bars give the evaporation, transpiration, irrigation, drainage and nitrogen leached of each.

\appfig{v4_0}{0.96}{Mulched against bare soil: evaporation, drainage and transpiration of the two surfaces.}

The line gives the difference between the setting of the user and bare soil in each flux.

\appfig{v4_1}{0.96}{Your setting against bare soil: the saving of each flux.}

The drawing is the picture of the loops of Chapter 9: rain and temperature, irrigation, the soil surface and its mulch, the crop, the drainage and the nitrogen, with the actual fluxes of the run (irrigation 250 mm, evaporation 172 mm, drainage 21 mm).

\appfig{v5_0}{0.96}{Feedback loops between climate, management, soil water, crop and nitrogen: the red arrows join the parts of the field that act on each other.}

Each row gives the effect of one step of one control (a $+0.5\,^\circ$C change, 10 per cent more rain, a stronger mulch) on the yield, the irrigation and the leaching.

\appfig{v5_1}{0.96}{What one step of each control does (click a row or a numbered arrow).}

Each panel is a question answered by sections and graphs of the model; for example the wetting and drying of the pores or the reach of the water table into the root zone.

\appfig{v5_2}{0.96}{Hidden questions of the vadose zone, answered for the coarse loam: sensors, compaction, the water table (questions 1 to 4).}

The last two panels connect the trigger to the pores that are emptied at that suction and the time that the soil delivers water to the surface.

\appfig{v5_3}{0.96}{Hidden questions of the vadose zone (questions 7 and 8): where on the suction scale the trigger really falls, and how long the soil keeps up with the sun.}

The figure divides the water that the mulch saves into less irrigation, more drainage, more transpiration and storage.

\appfig{n2_0}{0.96}{Fate of the water saved by mulch.}

The numbers are the fate of the saved millimetres: they are not all kept for the crop, and the panel follows the shares.

\appfig{n2_1}{0.96}{Fate of saved water: less irrigation 38 mm, more drainage 13 mm, more transpiration 17 mm.}

The five soil columns are the profile for five rainfall scales; the arrows show the evaporation and drainage and the colour bar gives the index that classifies the regime as drainage-dominated or evaporation-dominated.

\appfig{n3_0}{0.96}{Regime in the soil profile: vapour leaves upward, drainage downward, for rainfall scales from dry to wet.}

}{}
\IfFileExists{extra/long09.tex}{\subsection{Complete calculation of a mulch benefit}
\exbox{Mulch, tillage and the feedback loop}{Compare the evaporation under bare soil, straw mulch (28 per cent reduction) and plastic film (55 per cent) over a 10-day spell, find the length of stage-1 drying, the change of $K_s$ from tillage, and the feedback on the irrigation interval.}\sol
\ans Evaporation from bare soil: potential $E_0=4.5$ mm\,d$^{-1}$, mulch reduction factor $0.00$, so $E=E_0(1-k)=4.50$ mm\,d$^{-1}$; over a 10-day drying spell $45.0$ mm, saving $0.0$ mm
\ans Evaporation from straw mulch 5 t/ha: potential $E_0=4.5$ mm\,d$^{-1}$, mulch reduction factor $0.28$, so $E=E_0(1-k)=3.24$ mm\,d$^{-1}$; over a 10-day drying spell $32.4$ mm, saving $12.6$ mm
\ans Evaporation from plastic film: potential $E_0=4.5$ mm\,d$^{-1}$, mulch reduction factor $0.55$, so $E=E_0(1-k)=2.02$ mm\,d$^{-1}$; over a 10-day drying spell $20.2$ mm, saving $24.8$ mm
\ans Stage-1 evaporation (Ritchie) ends when $U=REW=9$ mm; with $E=4.5$ it lasts $9/4.5=2.0$ d on bare soil and $9/3.24=2.8$ d under straw
\ans Tillage: raising the porosity from 0.43 to 0.49 raises $K_s$ by the Kozeny--Carman ratio \[\left[\frac{0.49^3/(1-0.49)^2}{0.43^3/(1-0.43)^2}\right]=1.85\]
\ans Feedback loop: lower evaporation $\to$ wetter surface $\to$ longer interval; with a mulch saving of $12.6$ mm per 10 days the interval of 7 days becomes 8.3 days for the same TAW \finans


}{}
\IfFileExists{extra/proj09.tex}{\banner{Design project 9: a mulching and tillage plan for a tomato field}\label{sec:proj9}
\textbf{Brief.} A drip-irrigated tomato field on a loam (reference $ET_0=5$ mm\,d$^{-1}$, $a=3.5$ mm\,d$^{-1/2}$) is to be planted after four choices of surface cover. The designer compares the evaporation in the first thirty days after a wetting and the water saved. The evaporation after a wetting is $E=U+a(\sqrt{t}-\sqrt{t_1})$ with $t_1=U/ET_0$; with a cover the value is reduced by $(1-f_c)$.

\begin{table}[!htbp]\centering\small\caption{Thirty-day evaporation from a wetted loam for four surface covers ($ET_0=5$ mm/d, $a=3.5$ mm\,d$^{-1/2}$).}
\begin{tabular}{lrrrrr}\toprule
\hd{Cover} & \hd{Reduction factor} & \hd{$U$ (mm)} & \hd{$E_{30}$ (mm)} & \hd{Saving (mm)} & \hd{Saving (\%)}\\\midrule
Bare soil & 0.00 & 9 & 23.5 & 0.0 & 0\\
Straw 3 t/ha & 0.40 & 8 & 13.6 & 9.8 & 42\\
Straw 6 t/ha & 0.60 & 7 & 8.8 & 14.7 & 62\\
Plastic film & 0.90 & 5 & 2.1 & 21.4 & 91\\\bottomrule\end{tabular}\end{table}

\textbf{Reading the table.} Straw at 6 t\,ha$^{-1}$ saves 15 mm in thirty days, the plastic film 21 mm; over a 110-day season at the same rate, plastic could save as much as 78 mm, equal to the irrigation of three to four rounds. The choice depends on cost, on the rain (film sheds it) and on the disposal of the film.
\fig{g9_evap}{0.7}{Cumulative evaporation after a wetting for four surface covers ($ET_0=5$ mm\,d$^{-1}$, $a=3.5$ mm\,d$^{-1/2}$).}
}{}
\IfFileExists{extra/ten09.tex}{\subsection{Worked solutions on the ten soils}
Thirty days of evaporation after a wetting are computed for each soil, bare and under straw.

\exbox{Ten-soil example 9a: Coarse loam (S01)}{Coarse loam (class B: $U=9$ mm, $a=3.5$ mm\,d$^{-1/2}$ assumed): find the stage-one duration and the thirty-day evaporation after a wetting, bare and under straw at 6 t\,ha$^{-1}$ ($f_c=0.6$), at $ET_0=5.0$ mm\,d$^{-1}$.}\sol \givenline{$U=9$ mm ; $a=3.5$ mm\,d ; $f_c=0.6$ ; $ET_0=5.0$ mm\,d$^{-1}$}\ans $t_1=U/ET_0=9/5.0=1.8$ d \ans Bare: $E_{30}=U+a(\sqrt{30}-\sqrt{t_1})=9+3.5(5.477-1.342)=23.5$ mm \ans Under straw: $(1-0.6)\times23.5=9.4$ mm \ans Saving 14.1 mm (60 per cent) \finans

\exbox{Ten-soil example 9a: Coarse sand (S02)}{Coarse sand (class A: $U=6$ mm, $a=3.0$ mm\,d$^{-1/2}$ assumed): find the stage-one duration and the thirty-day evaporation after a wetting, bare and under straw at 6 t\,ha$^{-1}$ ($f_c=0.6$), at $ET_0=5.0$ mm\,d$^{-1}$.}\sol \givenline{$U=6$ mm ; $a=3.0$ mm\,d ; $f_c=0.6$ ; $ET_0=5.0$ mm\,d$^{-1}$}\ans $t_1=U/ET_0=6/5.0=1.2$ d \ans Bare: $E_{30}=U+a(\sqrt{30}-\sqrt{t_1})=6+3.0(5.477-1.095)=19.1$ mm \ans Under straw: $(1-0.6)\times19.1=7.7$ mm \ans Saving 11.5 mm (60 per cent) \finans

\exbox{Ten-soil example 9a: Sand (S03)}{Sand (class A: $U=6$ mm, $a=3.0$ mm\,d$^{-1/2}$ assumed): find the stage-one duration and the thirty-day evaporation after a wetting, bare and under straw at 6 t\,ha$^{-1}$ ($f_c=0.6$), at $ET_0=5.0$ mm\,d$^{-1}$.}\sol \givenline{$U=6$ mm ; $a=3.0$ mm\,d ; $f_c=0.6$ ; $ET_0=5.0$ mm\,d$^{-1}$}\ans $t_1=U/ET_0=6/5.0=1.2$ d \ans Bare: $E_{30}=U+a(\sqrt{30}-\sqrt{t_1})=6+3.0(5.477-1.095)=19.1$ mm \ans Under straw: $(1-0.6)\times19.1=7.7$ mm \ans Saving 11.5 mm (60 per cent) \finans

\exbox{Ten-soil example 9a: Loamy sand (S04)}{Loamy sand (class A: $U=6$ mm, $a=3.0$ mm\,d$^{-1/2}$ assumed): find the stage-one duration and the thirty-day evaporation after a wetting, bare and under straw at 6 t\,ha$^{-1}$ ($f_c=0.6$), at $ET_0=5.0$ mm\,d$^{-1}$.}\sol \givenline{$U=6$ mm ; $a=3.0$ mm\,d ; $f_c=0.6$ ; $ET_0=5.0$ mm\,d$^{-1}$}\ans $t_1=U/ET_0=6/5.0=1.2$ d \ans Bare: $E_{30}=U+a(\sqrt{30}-\sqrt{t_1})=6+3.0(5.477-1.095)=19.1$ mm \ans Under straw: $(1-0.6)\times19.1=7.7$ mm \ans Saving 11.5 mm (60 per cent) \finans

\exbox{Ten-soil example 9a: Sandy loam (S05)}{Sandy loam (class B: $U=9$ mm, $a=3.5$ mm\,d$^{-1/2}$ assumed): find the stage-one duration and the thirty-day evaporation after a wetting, bare and under straw at 6 t\,ha$^{-1}$ ($f_c=0.6$), at $ET_0=5.0$ mm\,d$^{-1}$.}\sol \givenline{$U=9$ mm ; $a=3.5$ mm\,d ; $f_c=0.6$ ; $ET_0=5.0$ mm\,d$^{-1}$}\ans $t_1=U/ET_0=9/5.0=1.8$ d \ans Bare: $E_{30}=U+a(\sqrt{30}-\sqrt{t_1})=9+3.5(5.477-1.342)=23.5$ mm \ans Under straw: $(1-0.6)\times23.5=9.4$ mm \ans Saving 14.1 mm (60 per cent) \finans

\exbox{Ten-soil example 9a: Loam (S06)}{Loam (class B: $U=9$ mm, $a=3.5$ mm\,d$^{-1/2}$ assumed): find the stage-one duration and the thirty-day evaporation after a wetting, bare and under straw at 6 t\,ha$^{-1}$ ($f_c=0.6$), at $ET_0=5.0$ mm\,d$^{-1}$.}\sol \givenline{$U=9$ mm ; $a=3.5$ mm\,d ; $f_c=0.6$ ; $ET_0=5.0$ mm\,d$^{-1}$}\ans $t_1=U/ET_0=9/5.0=1.8$ d \ans Bare: $E_{30}=U+a(\sqrt{30}-\sqrt{t_1})=9+3.5(5.477-1.342)=23.5$ mm \ans Under straw: $(1-0.6)\times23.5=9.4$ mm \ans Saving 14.1 mm (60 per cent) \finans

\exbox{Ten-soil example 9a: Silt loam (S07)}{Silt loam (class C: $U=10$ mm, $a=4.0$ mm\,d$^{-1/2}$ assumed): find the stage-one duration and the thirty-day evaporation after a wetting, bare and under straw at 6 t\,ha$^{-1}$ ($f_c=0.6$), at $ET_0=5.0$ mm\,d$^{-1}$.}\sol \givenline{$U=10$ mm ; $a=4.0$ mm\,d ; $f_c=0.6$ ; $ET_0=5.0$ mm\,d$^{-1}$}\ans $t_1=U/ET_0=10/5.0=2.0$ d \ans Bare: $E_{30}=U+a(\sqrt{30}-\sqrt{t_1})=10+4.0(5.477-1.414)=26.3$ mm \ans Under straw: $(1-0.6)\times26.3=10.5$ mm \ans Saving 15.8 mm (60 per cent) \finans

\exbox{Ten-soil example 9a: Sandy clay loam (S08)}{Sandy clay loam (class C: $U=10$ mm, $a=4.0$ mm\,d$^{-1/2}$ assumed): find the stage-one duration and the thirty-day evaporation after a wetting, bare and under straw at 6 t\,ha$^{-1}$ ($f_c=0.6$), at $ET_0=5.0$ mm\,d$^{-1}$.}\sol \givenline{$U=10$ mm ; $a=4.0$ mm\,d ; $f_c=0.6$ ; $ET_0=5.0$ mm\,d$^{-1}$}\ans $t_1=U/ET_0=10/5.0=2.0$ d \ans Bare: $E_{30}=U+a(\sqrt{30}-\sqrt{t_1})=10+4.0(5.477-1.414)=26.3$ mm \ans Under straw: $(1-0.6)\times26.3=10.5$ mm \ans Saving 15.8 mm (60 per cent) \finans

\exbox{Ten-soil example 9a: Clay loam (S09)}{Clay loam (class D: $U=12$ mm, $a=4.5$ mm\,d$^{-1/2}$ assumed): find the stage-one duration and the thirty-day evaporation after a wetting, bare and under straw at 6 t\,ha$^{-1}$ ($f_c=0.6$), at $ET_0=5.0$ mm\,d$^{-1}$.}\sol \givenline{$U=12$ mm ; $a=4.5$ mm\,d ; $f_c=0.6$ ; $ET_0=5.0$ mm\,d$^{-1}$}\ans $t_1=U/ET_0=12/5.0=2.4$ d \ans Bare: $E_{30}=U+a(\sqrt{30}-\sqrt{t_1})=12+4.5(5.477-1.549)=29.7$ mm \ans Under straw: $(1-0.6)\times29.7=11.9$ mm \ans Saving 17.8 mm (60 per cent) \finans

\exbox{Ten-soil example 9a: Silty clay loam (S10)}{Silty clay loam (class D: $U=12$ mm, $a=4.5$ mm\,d$^{-1/2}$ assumed): find the stage-one duration and the thirty-day evaporation after a wetting, bare and under straw at 6 t\,ha$^{-1}$ ($f_c=0.6$), at $ET_0=5.0$ mm\,d$^{-1}$.}\sol \givenline{$U=12$ mm ; $a=4.5$ mm\,d ; $f_c=0.6$ ; $ET_0=5.0$ mm\,d$^{-1}$}\ans $t_1=U/ET_0=12/5.0=2.4$ d \ans Bare: $E_{30}=U+a(\sqrt{30}-\sqrt{t_1})=12+4.5(5.477-1.549)=29.7$ mm \ans Under straw: $(1-0.6)\times29.7=11.9$ mm \ans Saving 17.8 mm (60 per cent) \finans

}{}
\IfFileExists{extra/tenb09.tex}{\subsection{Worked solutions on ten crops}
The ten crops of the library differ in their demand, their roots and their sensitivity. The same type of question is answered for each of them at $ET_0=5.5$ mm\,d$^{-1}$; the data are typical values of the crops (Doorenbos--Kassam and FAO 56).

\exbox{Crop example 9a: Maize}{Maize: split $K_c=1.20$ into basal $K_{cb}=0.95$ and an evaporation part $K_e=0.25$ of a bare wet soil; a straw mulch (6 t\,ha$^{-1}$) reduces $K_e$ by 60 per cent. Find the daily saving and the saving over 125 days if the soil is wet one day in three.}\sol \givenline{$K_c=1.20$ into basal ; $K_{cb}=0.95$ ; $K_e=0.25$}\ans $K_e=0.25$; under mulch $0.40\times0.25=0.10$ \ans Daily saving \[=0.60\times0.25\times5.5=0.82\ \text{mm\,d$^{-1}$}\] \ans Over the season: $0.82\times125/3=34$ mm \finans

\exbox{Crop example 9a: Wheat}{Wheat: split $K_c=1.15$ into basal $K_{cb}=0.90$ and an evaporation part $K_e=0.25$ of a bare wet soil; a straw mulch (6 t\,ha$^{-1}$) reduces $K_e$ by 60 per cent. Find the daily saving and the saving over 150 days if the soil is wet one day in three.}\sol \givenline{$K_c=1.15$ into basal ; $K_{cb}=0.90$ ; $K_e=0.25$}\ans $K_e=0.25$; under mulch $0.40\times0.25=0.10$ \ans Daily saving \[=0.60\times0.25\times5.5=0.82\ \text{mm\,d$^{-1}$}\] \ans Over the season: $0.82\times150/3=41$ mm \finans

\exbox{Crop example 9a: Cotton}{Cotton: split $K_c=1.15$ into basal $K_{cb}=0.90$ and an evaporation part $K_e=0.25$ of a bare wet soil; a straw mulch (6 t\,ha$^{-1}$) reduces $K_e$ by 60 per cent. Find the daily saving and the saving over 180 days if the soil is wet one day in three.}\sol \givenline{$K_c=1.15$ into basal ; $K_{cb}=0.90$ ; $K_e=0.25$}\ans $K_e=0.25$; under mulch $0.40\times0.25=0.10$ \ans Daily saving \[=0.60\times0.25\times5.5=0.82\ \text{mm\,d$^{-1}$}\] \ans Over the season: $0.82\times180/3=50$ mm \finans

\exbox{Crop example 9a: Tomato}{Tomato: split $K_c=1.15$ into basal $K_{cb}=0.90$ and an evaporation part $K_e=0.25$ of a bare wet soil; a straw mulch (6 t\,ha$^{-1}$) reduces $K_e$ by 60 per cent. Find the daily saving and the saving over 135 days if the soil is wet one day in three.}\sol \givenline{$K_c=1.15$ into basal ; $K_{cb}=0.90$ ; $K_e=0.25$}\ans $K_e=0.25$; under mulch $0.40\times0.25=0.10$ \ans Daily saving \[=0.60\times0.25\times5.5=0.82\ \text{mm\,d$^{-1}$}\] \ans Over the season: $0.82\times135/3=37$ mm \finans

\exbox{Crop example 9a: Potato}{Potato: split $K_c=1.15$ into basal $K_{cb}=0.90$ and an evaporation part $K_e=0.25$ of a bare wet soil; a straw mulch (6 t\,ha$^{-1}$) reduces $K_e$ by 60 per cent. Find the daily saving and the saving over 130 days if the soil is wet one day in three.}\sol \givenline{$K_c=1.15$ into basal ; $K_{cb}=0.90$ ; $K_e=0.25$}\ans $K_e=0.25$; under mulch $0.40\times0.25=0.10$ \ans Daily saving \[=0.60\times0.25\times5.5=0.82\ \text{mm\,d$^{-1}$}\] \ans Over the season: $0.82\times130/3=36$ mm \finans

\exbox{Crop example 9a: Soybean}{Soybean: split $K_c=1.15$ into basal $K_{cb}=0.90$ and an evaporation part $K_e=0.25$ of a bare wet soil; a straw mulch (6 t\,ha$^{-1}$) reduces $K_e$ by 60 per cent. Find the daily saving and the saving over 130 days if the soil is wet one day in three.}\sol \givenline{$K_c=1.15$ into basal ; $K_{cb}=0.90$ ; $K_e=0.25$}\ans $K_e=0.25$; under mulch $0.40\times0.25=0.10$ \ans Daily saving \[=0.60\times0.25\times5.5=0.82\ \text{mm\,d$^{-1}$}\] \ans Over the season: $0.82\times130/3=36$ mm \finans

\exbox{Crop example 9a: Sugar beet}{Sugar beet: split $K_c=1.20$ into basal $K_{cb}=0.95$ and an evaporation part $K_e=0.25$ of a bare wet soil; a straw mulch (6 t\,ha$^{-1}$) reduces $K_e$ by 60 per cent. Find the daily saving and the saving over 180 days if the soil is wet one day in three.}\sol \givenline{$K_c=1.20$ into basal ; $K_{cb}=0.95$ ; $K_e=0.25$}\ans $K_e=0.25$; under mulch $0.40\times0.25=0.10$ \ans Daily saving \[=0.60\times0.25\times5.5=0.82\ \text{mm\,d$^{-1}$}\] \ans Over the season: $0.82\times180/3=50$ mm \finans

\exbox{Crop example 9a: Sunflower}{Sunflower: split $K_c=1.10$ into basal $K_{cb}=0.85$ and an evaporation part $K_e=0.25$ of a bare wet soil; a straw mulch (6 t\,ha$^{-1}$) reduces $K_e$ by 60 per cent. Find the daily saving and the saving over 125 days if the soil is wet one day in three.}\sol \givenline{$K_c=1.10$ into basal ; $K_{cb}=0.85$ ; $K_e=0.25$}\ans $K_e=0.25$; under mulch $0.40\times0.25=0.10$ \ans Daily saving \[=0.60\times0.25\times5.5=0.82\ \text{mm\,d$^{-1}$}\] \ans Over the season: $0.82\times125/3=34$ mm \finans

\exbox{Crop example 9a: Onion}{Onion: split $K_c=1.05$ into basal $K_{cb}=0.80$ and an evaporation part $K_e=0.25$ of a bare wet soil; a straw mulch (6 t\,ha$^{-1}$) reduces $K_e$ by 60 per cent. Find the daily saving and the saving over 150 days if the soil is wet one day in three.}\sol \givenline{$K_c=1.05$ into basal ; $K_{cb}=0.80$ ; $K_e=0.25$}\ans $K_e=0.25$; under mulch $0.40\times0.25=0.10$ \ans Daily saving \[=0.60\times0.25\times5.5=0.82\ \text{mm\,d$^{-1}$}\] \ans Over the season: $0.82\times150/3=41$ mm \finans

\exbox{Crop example 9a: Alfalfa}{Alfalfa: split $K_c=1.15$ into basal $K_{cb}=0.90$ and an evaporation part $K_e=0.25$ of a bare wet soil; a straw mulch (6 t\,ha$^{-1}$) reduces $K_e$ by 60 per cent. Find the daily saving and the saving over 100 days if the soil is wet one day in three.}\sol \givenline{$K_c=1.15$ into basal ; $K_{cb}=0.90$ ; $K_e=0.25$}\ans $K_e=0.25$; under mulch $0.40\times0.25=0.10$ \ans Daily saving \[=0.60\times0.25\times5.5=0.82\ \text{mm\,d$^{-1}$}\] \ans Over the season: $0.82\times100/3=28$ mm \finans

}{}
\IfFileExists{extra/teny09.tex}{\subsection{Worked solutions on the ten soils (third scenario)}
The same questions are repeated for different working conditions, to show how the answer moves with them.

\exbox{Ten-soil example 9c: Coarse loam (S01)}{Coarse loam (class B: $U=9$ mm, $a=3.5$ mm\,d$^{-1/2}$ assumed): find the stage-one duration and the thirty-day evaporation after a wetting, bare and under straw at 6 t\,ha$^{-1}$ ($f_c=0.6$), at $ET_0=7.0$ mm\,d$^{-1}$.}\sol \givenline{$U=9$ mm ; $a=3.5$ mm\,d ; $f_c=0.6$ ; $ET_0=7.0$ mm\,d$^{-1}$}\ans $t_1=U/ET_0=9/7.0=1.3$ d \ans Bare: $E_{30}=U+a(\sqrt{30}-\sqrt{t_1})=9+3.5(5.477-1.134)=24.2$ mm \ans Under straw: $(1-0.6)\times24.2=9.7$ mm \ans Saving 14.5 mm (60 per cent) \finans

\exbox{Ten-soil example 9c: Coarse sand (S02)}{Coarse sand (class A: $U=6$ mm, $a=3.0$ mm\,d$^{-1/2}$ assumed): find the stage-one duration and the thirty-day evaporation after a wetting, bare and under straw at 6 t\,ha$^{-1}$ ($f_c=0.6$), at $ET_0=7.0$ mm\,d$^{-1}$.}\sol \givenline{$U=6$ mm ; $a=3.0$ mm\,d ; $f_c=0.6$ ; $ET_0=7.0$ mm\,d$^{-1}$}\ans $t_1=U/ET_0=6/7.0=0.9$ d \ans Bare: $E_{30}=U+a(\sqrt{30}-\sqrt{t_1})=6+3.0(5.477-0.926)=19.7$ mm \ans Under straw: $(1-0.6)\times19.7=7.9$ mm \ans Saving 11.8 mm (60 per cent) \finans

\exbox{Ten-soil example 9c: Sand (S03)}{Sand (class A: $U=6$ mm, $a=3.0$ mm\,d$^{-1/2}$ assumed): find the stage-one duration and the thirty-day evaporation after a wetting, bare and under straw at 6 t\,ha$^{-1}$ ($f_c=0.6$), at $ET_0=7.0$ mm\,d$^{-1}$.}\sol \givenline{$U=6$ mm ; $a=3.0$ mm\,d ; $f_c=0.6$ ; $ET_0=7.0$ mm\,d$^{-1}$}\ans $t_1=U/ET_0=6/7.0=0.9$ d \ans Bare: $E_{30}=U+a(\sqrt{30}-\sqrt{t_1})=6+3.0(5.477-0.926)=19.7$ mm \ans Under straw: $(1-0.6)\times19.7=7.9$ mm \ans Saving 11.8 mm (60 per cent) \finans

\exbox{Ten-soil example 9c: Loamy sand (S04)}{Loamy sand (class A: $U=6$ mm, $a=3.0$ mm\,d$^{-1/2}$ assumed): find the stage-one duration and the thirty-day evaporation after a wetting, bare and under straw at 6 t\,ha$^{-1}$ ($f_c=0.6$), at $ET_0=7.0$ mm\,d$^{-1}$.}\sol \givenline{$U=6$ mm ; $a=3.0$ mm\,d ; $f_c=0.6$ ; $ET_0=7.0$ mm\,d$^{-1}$}\ans $t_1=U/ET_0=6/7.0=0.9$ d \ans Bare: $E_{30}=U+a(\sqrt{30}-\sqrt{t_1})=6+3.0(5.477-0.926)=19.7$ mm \ans Under straw: $(1-0.6)\times19.7=7.9$ mm \ans Saving 11.8 mm (60 per cent) \finans

\exbox{Ten-soil example 9c: Sandy loam (S05)}{Sandy loam (class B: $U=9$ mm, $a=3.5$ mm\,d$^{-1/2}$ assumed): find the stage-one duration and the thirty-day evaporation after a wetting, bare and under straw at 6 t\,ha$^{-1}$ ($f_c=0.6$), at $ET_0=7.0$ mm\,d$^{-1}$.}\sol \givenline{$U=9$ mm ; $a=3.5$ mm\,d ; $f_c=0.6$ ; $ET_0=7.0$ mm\,d$^{-1}$}\ans $t_1=U/ET_0=9/7.0=1.3$ d \ans Bare: $E_{30}=U+a(\sqrt{30}-\sqrt{t_1})=9+3.5(5.477-1.134)=24.2$ mm \ans Under straw: $(1-0.6)\times24.2=9.7$ mm \ans Saving 14.5 mm (60 per cent) \finans

\exbox{Ten-soil example 9c: Loam (S06)}{Loam (class B: $U=9$ mm, $a=3.5$ mm\,d$^{-1/2}$ assumed): find the stage-one duration and the thirty-day evaporation after a wetting, bare and under straw at 6 t\,ha$^{-1}$ ($f_c=0.6$), at $ET_0=7.0$ mm\,d$^{-1}$.}\sol \givenline{$U=9$ mm ; $a=3.5$ mm\,d ; $f_c=0.6$ ; $ET_0=7.0$ mm\,d$^{-1}$}\ans $t_1=U/ET_0=9/7.0=1.3$ d \ans Bare: $E_{30}=U+a(\sqrt{30}-\sqrt{t_1})=9+3.5(5.477-1.134)=24.2$ mm \ans Under straw: $(1-0.6)\times24.2=9.7$ mm \ans Saving 14.5 mm (60 per cent) \finans

\exbox{Ten-soil example 9c: Silt loam (S07)}{Silt loam (class C: $U=10$ mm, $a=4.0$ mm\,d$^{-1/2}$ assumed): find the stage-one duration and the thirty-day evaporation after a wetting, bare and under straw at 6 t\,ha$^{-1}$ ($f_c=0.6$), at $ET_0=7.0$ mm\,d$^{-1}$.}\sol \givenline{$U=10$ mm ; $a=4.0$ mm\,d ; $f_c=0.6$ ; $ET_0=7.0$ mm\,d$^{-1}$}\ans $t_1=U/ET_0=10/7.0=1.4$ d \ans Bare: $E_{30}=U+a(\sqrt{30}-\sqrt{t_1})=10+4.0(5.477-1.195)=27.1$ mm \ans Under straw: $(1-0.6)\times27.1=10.9$ mm \ans Saving 16.3 mm (60 per cent) \finans

\exbox{Ten-soil example 9c: Sandy clay loam (S08)}{Sandy clay loam (class C: $U=10$ mm, $a=4.0$ mm\,d$^{-1/2}$ assumed): find the stage-one duration and the thirty-day evaporation after a wetting, bare and under straw at 6 t\,ha$^{-1}$ ($f_c=0.6$), at $ET_0=7.0$ mm\,d$^{-1}$.}\sol \givenline{$U=10$ mm ; $a=4.0$ mm\,d ; $f_c=0.6$ ; $ET_0=7.0$ mm\,d$^{-1}$}\ans $t_1=U/ET_0=10/7.0=1.4$ d \ans Bare: $E_{30}=U+a(\sqrt{30}-\sqrt{t_1})=10+4.0(5.477-1.195)=27.1$ mm \ans Under straw: $(1-0.6)\times27.1=10.9$ mm \ans Saving 16.3 mm (60 per cent) \finans

\exbox{Ten-soil example 9c: Clay loam (S09)}{Clay loam (class D: $U=12$ mm, $a=4.5$ mm\,d$^{-1/2}$ assumed): find the stage-one duration and the thirty-day evaporation after a wetting, bare and under straw at 6 t\,ha$^{-1}$ ($f_c=0.6$), at $ET_0=7.0$ mm\,d$^{-1}$.}\sol \givenline{$U=12$ mm ; $a=4.5$ mm\,d ; $f_c=0.6$ ; $ET_0=7.0$ mm\,d$^{-1}$}\ans $t_1=U/ET_0=12/7.0=1.7$ d \ans Bare: $E_{30}=U+a(\sqrt{30}-\sqrt{t_1})=12+4.5(5.477-1.309)=30.8$ mm \ans Under straw: $(1-0.6)\times30.8=12.3$ mm \ans Saving 18.5 mm (60 per cent) \finans

\exbox{Ten-soil example 9c: Silty clay loam (S10)}{Silty clay loam (class D: $U=12$ mm, $a=4.5$ mm\,d$^{-1/2}$ assumed): find the stage-one duration and the thirty-day evaporation after a wetting, bare and under straw at 6 t\,ha$^{-1}$ ($f_c=0.6$), at $ET_0=7.0$ mm\,d$^{-1}$.}\sol \givenline{$U=12$ mm ; $a=4.5$ mm\,d ; $f_c=0.6$ ; $ET_0=7.0$ mm\,d$^{-1}$}\ans $t_1=U/ET_0=12/7.0=1.7$ d \ans Bare: $E_{30}=U+a(\sqrt{30}-\sqrt{t_1})=12+4.5(5.477-1.309)=30.8$ mm \ans Under straw: $(1-0.6)\times30.8=12.3$ mm \ans Saving 18.5 mm (60 per cent) \finans

}{}
\IfFileExists{extra/tenby09.tex}{\subsection{Worked solutions on ten crops (hot climate)}
The ten crops of the library differ in their demand, their roots and their sensitivity. The same type of question is answered for each of them at $ET_0=7.5$ mm\,d$^{-1}$; the data are typical values of the crops (Doorenbos--Kassam and FAO 56).

\exbox{Crop example 9c: Maize}{Maize: split $K_c=1.20$ into basal $K_{cb}=0.95$ and an evaporation part $K_e=0.25$ of a bare wet soil; a straw mulch (6 t\,ha$^{-1}$) reduces $K_e$ by 60 per cent. Find the daily saving and the saving over 125 days if the soil is wet one day in three.}\sol $K_e=0.25$; under mulch $0.40\times0.25=0.10$ \ans Daily saving $=0.60\times0.25\times7.5=1.12$ mm\,d$^{-1}$ \ans Over the season: $1.12\times125/3=47$ mm \ans

\exbox{Crop example 9c: Wheat}{Wheat: split $K_c=1.15$ into basal $K_{cb}=0.90$ and an evaporation part $K_e=0.25$ of a bare wet soil; a straw mulch (6 t\,ha$^{-1}$) reduces $K_e$ by 60 per cent. Find the daily saving and the saving over 150 days if the soil is wet one day in three.}\sol $K_e=0.25$; under mulch $0.40\times0.25=0.10$ \ans Daily saving $=0.60\times0.25\times7.5=1.12$ mm\,d$^{-1}$ \ans Over the season: $1.12\times150/3=56$ mm \ans

\exbox{Crop example 9c: Cotton}{Cotton: split $K_c=1.15$ into basal $K_{cb}=0.90$ and an evaporation part $K_e=0.25$ of a bare wet soil; a straw mulch (6 t\,ha$^{-1}$) reduces $K_e$ by 60 per cent. Find the daily saving and the saving over 180 days if the soil is wet one day in three.}\sol $K_e=0.25$; under mulch $0.40\times0.25=0.10$ \ans Daily saving $=0.60\times0.25\times7.5=1.12$ mm\,d$^{-1}$ \ans Over the season: $1.12\times180/3=68$ mm \ans

\exbox{Crop example 9c: Tomato}{Tomato: split $K_c=1.15$ into basal $K_{cb}=0.90$ and an evaporation part $K_e=0.25$ of a bare wet soil; a straw mulch (6 t\,ha$^{-1}$) reduces $K_e$ by 60 per cent. Find the daily saving and the saving over 135 days if the soil is wet one day in three.}\sol $K_e=0.25$; under mulch $0.40\times0.25=0.10$ \ans Daily saving $=0.60\times0.25\times7.5=1.12$ mm\,d$^{-1}$ \ans Over the season: $1.12\times135/3=51$ mm \ans

\exbox{Crop example 9c: Potato}{Potato: split $K_c=1.15$ into basal $K_{cb}=0.90$ and an evaporation part $K_e=0.25$ of a bare wet soil; a straw mulch (6 t\,ha$^{-1}$) reduces $K_e$ by 60 per cent. Find the daily saving and the saving over 130 days if the soil is wet one day in three.}\sol $K_e=0.25$; under mulch $0.40\times0.25=0.10$ \ans Daily saving $=0.60\times0.25\times7.5=1.12$ mm\,d$^{-1}$ \ans Over the season: $1.12\times130/3=49$ mm \ans

\exbox{Crop example 9c: Soybean}{Soybean: split $K_c=1.15$ into basal $K_{cb}=0.90$ and an evaporation part $K_e=0.25$ of a bare wet soil; a straw mulch (6 t\,ha$^{-1}$) reduces $K_e$ by 60 per cent. Find the daily saving and the saving over 130 days if the soil is wet one day in three.}\sol $K_e=0.25$; under mulch $0.40\times0.25=0.10$ \ans Daily saving $=0.60\times0.25\times7.5=1.12$ mm\,d$^{-1}$ \ans Over the season: $1.12\times130/3=49$ mm \ans

\exbox{Crop example 9c: Sugar beet}{Sugar beet: split $K_c=1.20$ into basal $K_{cb}=0.95$ and an evaporation part $K_e=0.25$ of a bare wet soil; a straw mulch (6 t\,ha$^{-1}$) reduces $K_e$ by 60 per cent. Find the daily saving and the saving over 180 days if the soil is wet one day in three.}\sol $K_e=0.25$; under mulch $0.40\times0.25=0.10$ \ans Daily saving $=0.60\times0.25\times7.5=1.12$ mm\,d$^{-1}$ \ans Over the season: $1.12\times180/3=68$ mm \ans

\exbox{Crop example 9c: Sunflower}{Sunflower: split $K_c=1.10$ into basal $K_{cb}=0.85$ and an evaporation part $K_e=0.25$ of a bare wet soil; a straw mulch (6 t\,ha$^{-1}$) reduces $K_e$ by 60 per cent. Find the daily saving and the saving over 125 days if the soil is wet one day in three.}\sol $K_e=0.25$; under mulch $0.40\times0.25=0.10$ \ans Daily saving $=0.60\times0.25\times7.5=1.12$ mm\,d$^{-1}$ \ans Over the season: $1.12\times125/3=47$ mm \ans

\exbox{Crop example 9c: Onion}{Onion: split $K_c=1.05$ into basal $K_{cb}=0.80$ and an evaporation part $K_e=0.25$ of a bare wet soil; a straw mulch (6 t\,ha$^{-1}$) reduces $K_e$ by 60 per cent. Find the daily saving and the saving over 150 days if the soil is wet one day in three.}\sol $K_e=0.25$; under mulch $0.40\times0.25=0.10$ \ans Daily saving $=0.60\times0.25\times7.5=1.12$ mm\,d$^{-1}$ \ans Over the season: $1.12\times150/3=56$ mm \ans

\exbox{Crop example 9c: Alfalfa}{Alfalfa: split $K_c=1.15$ into basal $K_{cb}=0.90$ and an evaporation part $K_e=0.25$ of a bare wet soil; a straw mulch (6 t\,ha$^{-1}$) reduces $K_e$ by 60 per cent. Find the daily saving and the saving over 100 days if the soil is wet one day in three.}\sol $K_e=0.25$; under mulch $0.40\times0.25=0.10$ \ans Daily saving $=0.60\times0.25\times7.5=1.12$ mm\,d$^{-1}$ \ans Over the season: $1.12\times100/3=38$ mm \ans

}{}
\IfFileExists{extra/tenx09.tex}{\subsection{Worked solutions on the ten soils (second scenario)}
The same questions are repeated for different working conditions, to show how the answer moves with them.

\exbox{Ten-soil example 9b: Coarse loam (S01)}{Coarse loam (class B: $U=9$ mm, $a=3.5$ mm\,d$^{-1/2}$ assumed): find the stage-one duration and the thirty-day evaporation after a wetting, bare and under straw at 6 t\,ha$^{-1}$ ($f_c=0.6$), at $ET_0=3.5$ mm\,d$^{-1}$.}\sol \givenline{$U=9$ mm ; $a=3.5$ mm\,d ; $f_c=0.6$ ; $ET_0=3.5$ mm\,d$^{-1}$}\ans $t_1=U/ET_0=9/3.5=2.6$ d \ans Bare: $E_{30}=U+a(\sqrt{30}-\sqrt{t_1})=9+3.5(5.477-1.604)=22.6$ mm \ans Under straw: $(1-0.6)\times22.6=9.0$ mm \ans Saving 13.5 mm (60 per cent) \finans

\exbox{Ten-soil example 9b: Coarse sand (S02)}{Coarse sand (class A: $U=6$ mm, $a=3.0$ mm\,d$^{-1/2}$ assumed): find the stage-one duration and the thirty-day evaporation after a wetting, bare and under straw at 6 t\,ha$^{-1}$ ($f_c=0.6$), at $ET_0=3.5$ mm\,d$^{-1}$.}\sol \givenline{$U=6$ mm ; $a=3.0$ mm\,d ; $f_c=0.6$ ; $ET_0=3.5$ mm\,d$^{-1}$}\ans $t_1=U/ET_0=6/3.5=1.7$ d \ans Bare: $E_{30}=U+a(\sqrt{30}-\sqrt{t_1})=6+3.0(5.477-1.309)=18.5$ mm \ans Under straw: $(1-0.6)\times18.5=7.4$ mm \ans Saving 11.1 mm (60 per cent) \finans

\exbox{Ten-soil example 9b: Sand (S03)}{Sand (class A: $U=6$ mm, $a=3.0$ mm\,d$^{-1/2}$ assumed): find the stage-one duration and the thirty-day evaporation after a wetting, bare and under straw at 6 t\,ha$^{-1}$ ($f_c=0.6$), at $ET_0=3.5$ mm\,d$^{-1}$.}\sol \givenline{$U=6$ mm ; $a=3.0$ mm\,d ; $f_c=0.6$ ; $ET_0=3.5$ mm\,d$^{-1}$}\ans $t_1=U/ET_0=6/3.5=1.7$ d \ans Bare: $E_{30}=U+a(\sqrt{30}-\sqrt{t_1})=6+3.0(5.477-1.309)=18.5$ mm \ans Under straw: $(1-0.6)\times18.5=7.4$ mm \ans Saving 11.1 mm (60 per cent) \finans

\exbox{Ten-soil example 9b: Loamy sand (S04)}{Loamy sand (class A: $U=6$ mm, $a=3.0$ mm\,d$^{-1/2}$ assumed): find the stage-one duration and the thirty-day evaporation after a wetting, bare and under straw at 6 t\,ha$^{-1}$ ($f_c=0.6$), at $ET_0=3.5$ mm\,d$^{-1}$.}\sol \givenline{$U=6$ mm ; $a=3.0$ mm\,d ; $f_c=0.6$ ; $ET_0=3.5$ mm\,d$^{-1}$}\ans $t_1=U/ET_0=6/3.5=1.7$ d \ans Bare: $E_{30}=U+a(\sqrt{30}-\sqrt{t_1})=6+3.0(5.477-1.309)=18.5$ mm \ans Under straw: $(1-0.6)\times18.5=7.4$ mm \ans Saving 11.1 mm (60 per cent) \finans

\exbox{Ten-soil example 9b: Sandy loam (S05)}{Sandy loam (class B: $U=9$ mm, $a=3.5$ mm\,d$^{-1/2}$ assumed): find the stage-one duration and the thirty-day evaporation after a wetting, bare and under straw at 6 t\,ha$^{-1}$ ($f_c=0.6$), at $ET_0=3.5$ mm\,d$^{-1}$.}\sol \givenline{$U=9$ mm ; $a=3.5$ mm\,d ; $f_c=0.6$ ; $ET_0=3.5$ mm\,d$^{-1}$}\ans $t_1=U/ET_0=9/3.5=2.6$ d \ans Bare: $E_{30}=U+a(\sqrt{30}-\sqrt{t_1})=9+3.5(5.477-1.604)=22.6$ mm \ans Under straw: $(1-0.6)\times22.6=9.0$ mm \ans Saving 13.5 mm (60 per cent) \finans

\exbox{Ten-soil example 9b: Loam (S06)}{Loam (class B: $U=9$ mm, $a=3.5$ mm\,d$^{-1/2}$ assumed): find the stage-one duration and the thirty-day evaporation after a wetting, bare and under straw at 6 t\,ha$^{-1}$ ($f_c=0.6$), at $ET_0=3.5$ mm\,d$^{-1}$.}\sol \givenline{$U=9$ mm ; $a=3.5$ mm\,d ; $f_c=0.6$ ; $ET_0=3.5$ mm\,d$^{-1}$}\ans $t_1=U/ET_0=9/3.5=2.6$ d \ans Bare: $E_{30}=U+a(\sqrt{30}-\sqrt{t_1})=9+3.5(5.477-1.604)=22.6$ mm \ans Under straw: $(1-0.6)\times22.6=9.0$ mm \ans Saving 13.5 mm (60 per cent) \finans

\exbox{Ten-soil example 9b: Silt loam (S07)}{Silt loam (class C: $U=10$ mm, $a=4.0$ mm\,d$^{-1/2}$ assumed): find the stage-one duration and the thirty-day evaporation after a wetting, bare and under straw at 6 t\,ha$^{-1}$ ($f_c=0.6$), at $ET_0=3.5$ mm\,d$^{-1}$.}\sol \givenline{$U=10$ mm ; $a=4.0$ mm\,d ; $f_c=0.6$ ; $ET_0=3.5$ mm\,d$^{-1}$}\ans $t_1=U/ET_0=10/3.5=2.9$ d \ans Bare: $E_{30}=U+a(\sqrt{30}-\sqrt{t_1})=10+4.0(5.477-1.690)=25.1$ mm \ans Under straw: $(1-0.6)\times25.1=10.1$ mm \ans Saving 15.1 mm (60 per cent) \finans

\exbox{Ten-soil example 9b: Sandy clay loam (S08)}{Sandy clay loam (class C: $U=10$ mm, $a=4.0$ mm\,d$^{-1/2}$ assumed): find the stage-one duration and the thirty-day evaporation after a wetting, bare and under straw at 6 t\,ha$^{-1}$ ($f_c=0.6$), at $ET_0=3.5$ mm\,d$^{-1}$.}\sol \givenline{$U=10$ mm ; $a=4.0$ mm\,d ; $f_c=0.6$ ; $ET_0=3.5$ mm\,d$^{-1}$}\ans $t_1=U/ET_0=10/3.5=2.9$ d \ans Bare: $E_{30}=U+a(\sqrt{30}-\sqrt{t_1})=10+4.0(5.477-1.690)=25.1$ mm \ans Under straw: $(1-0.6)\times25.1=10.1$ mm \ans Saving 15.1 mm (60 per cent) \finans

\exbox{Ten-soil example 9b: Clay loam (S09)}{Clay loam (class D: $U=12$ mm, $a=4.5$ mm\,d$^{-1/2}$ assumed): find the stage-one duration and the thirty-day evaporation after a wetting, bare and under straw at 6 t\,ha$^{-1}$ ($f_c=0.6$), at $ET_0=3.5$ mm\,d$^{-1}$.}\sol \givenline{$U=12$ mm ; $a=4.5$ mm\,d ; $f_c=0.6$ ; $ET_0=3.5$ mm\,d$^{-1}$}\ans $t_1=U/ET_0=12/3.5=3.4$ d \ans Bare: $E_{30}=U+a(\sqrt{30}-\sqrt{t_1})=12+4.5(5.477-1.852)=28.3$ mm \ans Under straw: $(1-0.6)\times28.3=11.3$ mm \ans Saving 17.0 mm (60 per cent) \finans

\exbox{Ten-soil example 9b: Silty clay loam (S10)}{Silty clay loam (class D: $U=12$ mm, $a=4.5$ mm\,d$^{-1/2}$ assumed): find the stage-one duration and the thirty-day evaporation after a wetting, bare and under straw at 6 t\,ha$^{-1}$ ($f_c=0.6$), at $ET_0=3.5$ mm\,d$^{-1}$.}\sol \givenline{$U=12$ mm ; $a=4.5$ mm\,d ; $f_c=0.6$ ; $ET_0=3.5$ mm\,d$^{-1}$}\ans $t_1=U/ET_0=12/3.5=3.4$ d \ans Bare: $E_{30}=U+a(\sqrt{30}-\sqrt{t_1})=12+4.5(5.477-1.852)=28.3$ mm \ans Under straw: $(1-0.6)\times28.3=11.3$ mm \ans Saving 17.0 mm (60 per cent) \finans

}{}
\IfFileExists{extra/tenbx09.tex}{\subsection{Worked solutions on ten crops (cool climate)}
The ten crops of the library differ in their demand, their roots and their sensitivity. The same type of question is answered for each of them at $ET_0=4.0$ mm\,d$^{-1}$; the data are typical values of the crops (Doorenbos--Kassam and FAO 56).

\exbox{Crop example 9b: Maize}{Maize: split $K_c=1.20$ into basal $K_{cb}=0.95$ and an evaporation part $K_e=0.25$ of a bare wet soil; a straw mulch (6 t\,ha$^{-1}$) reduces $K_e$ by 60 per cent. Find the daily saving and the saving over 125 days if the soil is wet one day in three.}\sol \givenline{$K_c=1.20$ into basal ; $K_{cb}=0.95$ ; $K_e=0.25$}\ans $K_e=0.25$; under mulch $0.40\times0.25=0.10$ \ans Daily saving $=0.60\times0.25\times4.0=0.60$ mm\,d$^{-1}$ \ans Over the season: $0.60\times125/3=25$ mm \finans

\exbox{Crop example 9b: Wheat}{Wheat: split $K_c=1.15$ into basal $K_{cb}=0.90$ and an evaporation part $K_e=0.25$ of a bare wet soil; a straw mulch (6 t\,ha$^{-1}$) reduces $K_e$ by 60 per cent. Find the daily saving and the saving over 150 days if the soil is wet one day in three.}\sol \givenline{$K_c=1.15$ into basal ; $K_{cb}=0.90$ ; $K_e=0.25$}\ans $K_e=0.25$; under mulch $0.40\times0.25=0.10$ \ans Daily saving $=0.60\times0.25\times4.0=0.60$ mm\,d$^{-1}$ \ans Over the season: $0.60\times150/3=30$ mm \finans

\exbox{Crop example 9b: Cotton}{Cotton: split $K_c=1.15$ into basal $K_{cb}=0.90$ and an evaporation part $K_e=0.25$ of a bare wet soil; a straw mulch (6 t\,ha$^{-1}$) reduces $K_e$ by 60 per cent. Find the daily saving and the saving over 180 days if the soil is wet one day in three.}\sol \givenline{$K_c=1.15$ into basal ; $K_{cb}=0.90$ ; $K_e=0.25$}\ans $K_e=0.25$; under mulch $0.40\times0.25=0.10$ \ans Daily saving $=0.60\times0.25\times4.0=0.60$ mm\,d$^{-1}$ \ans Over the season: $0.60\times180/3=36$ mm \finans

\exbox{Crop example 9b: Tomato}{Tomato: split $K_c=1.15$ into basal $K_{cb}=0.90$ and an evaporation part $K_e=0.25$ of a bare wet soil; a straw mulch (6 t\,ha$^{-1}$) reduces $K_e$ by 60 per cent. Find the daily saving and the saving over 135 days if the soil is wet one day in three.}\sol \givenline{$K_c=1.15$ into basal ; $K_{cb}=0.90$ ; $K_e=0.25$}\ans $K_e=0.25$; under mulch $0.40\times0.25=0.10$ \ans Daily saving $=0.60\times0.25\times4.0=0.60$ mm\,d$^{-1}$ \ans Over the season: $0.60\times135/3=27$ mm \finans

\exbox{Crop example 9b: Potato}{Potato: split $K_c=1.15$ into basal $K_{cb}=0.90$ and an evaporation part $K_e=0.25$ of a bare wet soil; a straw mulch (6 t\,ha$^{-1}$) reduces $K_e$ by 60 per cent. Find the daily saving and the saving over 130 days if the soil is wet one day in three.}\sol \givenline{$K_c=1.15$ into basal ; $K_{cb}=0.90$ ; $K_e=0.25$}\ans $K_e=0.25$; under mulch $0.40\times0.25=0.10$ \ans Daily saving $=0.60\times0.25\times4.0=0.60$ mm\,d$^{-1}$ \ans Over the season: $0.60\times130/3=26$ mm \finans

\exbox{Crop example 9b: Soybean}{Soybean: split $K_c=1.15$ into basal $K_{cb}=0.90$ and an evaporation part $K_e=0.25$ of a bare wet soil; a straw mulch (6 t\,ha$^{-1}$) reduces $K_e$ by 60 per cent. Find the daily saving and the saving over 130 days if the soil is wet one day in three.}\sol \givenline{$K_c=1.15$ into basal ; $K_{cb}=0.90$ ; $K_e=0.25$}\ans $K_e=0.25$; under mulch $0.40\times0.25=0.10$ \ans Daily saving $=0.60\times0.25\times4.0=0.60$ mm\,d$^{-1}$ \ans Over the season: $0.60\times130/3=26$ mm \finans

\exbox{Crop example 9b: Sugar beet}{Sugar beet: split $K_c=1.20$ into basal $K_{cb}=0.95$ and an evaporation part $K_e=0.25$ of a bare wet soil; a straw mulch (6 t\,ha$^{-1}$) reduces $K_e$ by 60 per cent. Find the daily saving and the saving over 180 days if the soil is wet one day in three.}\sol \givenline{$K_c=1.20$ into basal ; $K_{cb}=0.95$ ; $K_e=0.25$}\ans $K_e=0.25$; under mulch $0.40\times0.25=0.10$ \ans Daily saving $=0.60\times0.25\times4.0=0.60$ mm\,d$^{-1}$ \ans Over the season: $0.60\times180/3=36$ mm \finans

\exbox{Crop example 9b: Sunflower}{Sunflower: split $K_c=1.10$ into basal $K_{cb}=0.85$ and an evaporation part $K_e=0.25$ of a bare wet soil; a straw mulch (6 t\,ha$^{-1}$) reduces $K_e$ by 60 per cent. Find the daily saving and the saving over 125 days if the soil is wet one day in three.}\sol \givenline{$K_c=1.10$ into basal ; $K_{cb}=0.85$ ; $K_e=0.25$}\ans $K_e=0.25$; under mulch $0.40\times0.25=0.10$ \ans Daily saving $=0.60\times0.25\times4.0=0.60$ mm\,d$^{-1}$ \ans Over the season: $0.60\times125/3=25$ mm \finans

\exbox{Crop example 9b: Onion}{Onion: split $K_c=1.05$ into basal $K_{cb}=0.80$ and an evaporation part $K_e=0.25$ of a bare wet soil; a straw mulch (6 t\,ha$^{-1}$) reduces $K_e$ by 60 per cent. Find the daily saving and the saving over 150 days if the soil is wet one day in three.}\sol \givenline{$K_c=1.05$ into basal ; $K_{cb}=0.80$ ; $K_e=0.25$}\ans $K_e=0.25$; under mulch $0.40\times0.25=0.10$ \ans Daily saving $=0.60\times0.25\times4.0=0.60$ mm\,d$^{-1}$ \ans Over the season: $0.60\times150/3=30$ mm \finans

\exbox{Crop example 9b: Alfalfa}{Alfalfa: split $K_c=1.15$ into basal $K_{cb}=0.90$ and an evaporation part $K_e=0.25$ of a bare wet soil; a straw mulch (6 t\,ha$^{-1}$) reduces $K_e$ by 60 per cent. Find the daily saving and the saving over 100 days if the soil is wet one day in three.}\sol \givenline{$K_c=1.15$ into basal ; $K_{cb}=0.90$ ; $K_e=0.25$}\ans $K_e=0.25$; under mulch $0.40\times0.25=0.10$ \ans Daily saving $=0.60\times0.25\times4.0=0.60$ mm\,d$^{-1}$ \ans Over the season: $0.60\times100/3=20$ mm \finans

}{}
\IfFileExists{extra/ex09.tex}{\backhead{EXERCISES}{khred}
\begin{exlist}
\item Find $\phi$ for $\rho_b=1.35$.\exans{0.491}
\item A 20-cm layer at 1.50 is loosened to 1.25. Swelling?\exans{4.0 cm}
\item Find $S_d$ for $RR=3$ cm, $s=4$.\exans{0.47 mm}
\item Fixed point $K^*$ for $\alpha=0.03$, $K_{max}=30$, $\beta=0.0004$, $R=50$.\exans{18.0}
\item $E_{30}$ for $U=8$, $a=3$, $ET_0=4$ (no cover).\exans{20.2 mm}
\item The same with $f_c=0.5$.\exans{10.1 mm}
\item Name two reinforcing and two balancing feedbacks in the soil--crop system.\exans{see text}
\item State three effects of tillage on the surface balance.\exans{roughness; infiltration; evaporation supply}
\item Why is plastic film unsuitable with sprinklers?\exans{it sheds the water}
\item Describe the structure--infiltration loop.\exans{tillage up, slaking down}
\end{exlist}
}{}
\IfFileExists{extra/dyk09.tex}{\dykband{DO YOU KNOW ?}
\begin{exlist}
\item What is the difference between total porosity and macro-porosity?
\item How does tillage change the surface storage?
\item List three kinds of mulch and one effect of each.
\item Why does mulch shorten stage one?
\item What is a reinforcing feedback?
\item Write the fixed point of the structure--infiltration loop.
\item What does the cover factor $f_c$ do in the model?
\item Why must film be removed after the season?
\end{exlist}
}{}
\IfFileExists{extra/obj09.tex}{\dykband{OBJECTIVE TYPE QUESTIONS}
\begin{exlist}
\item $\phi$ is\opt{$1-\rho_b/\rho_s$}{$\rho_b/\rho_s$}{$\rho_s-\rho_b$}{$\rho_b\rho_s$}
\item No-till\opt{loosens the soil}{leaves residue on top}{buries residue}{stops rain}
\item Plastic film\opt{adds rain}{blocks evaporation}{lowers $\rho_b$}{feeds biology}
\item Stage 1 evaporation is limited by\opt{soil supply}{energy}{wind}{the crop}
\item A balancing loop\opt{amplifies}{opposes change}{has no state}{is random}
\item The roughness stores\opt{gas}{water in depressions}{salt}{heat}
\item The cover factor lies in\opt{[0,1]}{[1,2]}{[-1,0]}{any}
\item Residue decomposes\opt{never}{by biology}{only in winter}{by plastic}
\item Tillage halves\opt{always runoff}{nothing fixed}{gravity}{rain}
\item Organic mulch adds\opt{organic matter}{glass}{iron}{none}
\end{exlist}
\objkey{1.~(a)\quad 2.~(b)\quad 3.~(b)\quad 4.~(b)\quad 5.~(b)\quad 6.~(b)\quad 7.~(a)\quad 8.~(b)\quad 9.~(b)\quad 10.~(a)}
}{} \finans